\appendix
\ctexset{section={aftername={\quad}}}
\section{关键源代码节选}\label{app:source-code}
以下片段对应本文的四个模型步骤：计算图预处理、候选方案构造、轻量排序和评价接口，仅列出能够对应公式、约束和事件定义的关键实现。

\begingroup
\IfFontExistsTF{Menlo}{\setmonofont{Menlo}}{\setmonofont{Consolas}}
\fvset{fontsize=\scriptsize,baselinestretch=1,breaklines=true,breakanywhere=true,breakanywheresymbolpre={},breaksymbolleft={},numbers=left,numbersep=4pt,xleftmargin=2em}

\subsection{计算图与任务优先级}
\noindent\textbf{计算图与任务优先级：}张量依赖收缩、关键路径和释放优先级。\par\nobreak
\VerbatimInput[firstline=22,lastline=70,firstnumber=22]{support/experiment/src/huawei_code/graph.py}
\VerbatimInput[firstline=91,lastline=107,firstnumber=91]{support/experiment/src/huawei_code/graph.py}
\VerbatimInput[firstline=110,lastline=141,firstnumber=110]{support/experiment/src/huawei_code/graph.py}
\VerbatimInput[firstline=144,lastline=170,firstnumber=144]{support/experiment/src/huawei_code/graph.py}

\subsection{候选方案构造与约束}
\noindent\textbf{候选方案构造与约束：}空闲核继承、规范化编号、场景化依赖合法性检查和初始分块方案。\par\nobreak
\VerbatimInput[firstline=19,lastline=96,firstnumber=19]{support/experiment/src/huawei_code/plan.py}
\VerbatimInput[firstline=110,lastline=150,firstnumber=110]{support/experiment/src/huawei_code/plan.py}

\noindent\textbf{候选变换：}局部移动、聚合、重排及共享张量的Cache候选。\par\nobreak
\VerbatimInput[firstline=13,lastline=24,firstnumber=13]{support/experiment/src/huawei_code/candidates.py}
\VerbatimInput[firstline=27,lastline=78,firstnumber=27]{support/experiment/src/huawei_code/candidates.py}
\VerbatimInput[firstline=81,lastline=97,firstnumber=81]{support/experiment/src/huawei_code/candidates.py}
\VerbatimInput[firstline=100,lastline=140,firstnumber=100]{support/experiment/src/huawei_code/candidates.py}

\subsection{轻量排序与官方评价}
\noindent\textbf{轻量排序：}互斥结构搬运、容量压力、逻辑COPY查询代理及候选排序。\par\nobreak
\VerbatimInput[firstline=31,lastline=80,firstnumber=31]{support/experiment/src/huawei_code/scoring.py}
\VerbatimInput[firstline=83,lastline=119,firstnumber=83]{support/experiment/src/huawei_code/scoring.py}
\VerbatimInput[firstline=127,lastline=206,firstnumber=127]{support/experiment/src/huawei_code/scoring.py}

\noindent\textbf{评价接口：}调用题面评价程序并分别返回A、B、C三问的工期、搬运和命中字段。\par\nobreak
\VerbatimInput[firstline=11,lastline=16,firstnumber=11]{support/experiment/src/huawei_code/official.py}
\VerbatimInput[firstline=44,lastline=58,firstnumber=44]{support/experiment/src/huawei_code/official.py}
\endgroup

\section{数据与指标说明}\label{app:provenance}
表\ref{tab:a-code-map}列出与模型步骤对应的实现模块；表\ref{tab:b-artifacts}列出每个图--核数--场景保存的方案与评价记录，其中\texttt{profile.json}为展开检查结果，\texttt{result.json}为完整事件模拟结果。
\begin{table}[H]\centering\small\caption{实现模块与模型步骤}\label{tab:a-code-map}
\begin{tabularx}{\linewidth}{lYY}\toprule
实现模块 & 模型步骤 & 输出对象\\\midrule
\texttt{graph.py}、\texttt{plan.py} & 生产消费索引、D/H序、四起点与计划合法性检查 & 依赖索引、映射和核内序列\\
\texttt{candidates.py}、\texttt{scoring.py} & 局部变体、去重与轻量排序 & 候选方案和评分向量\\
\texttt{official.py}、题给评估程序 & Task展开、展开检查与共享DDR事件模拟 & Task/COPY展开、工期和搬运\\
\texttt{official.py}的选优接口 & 成功结果按工期、搬运排序 & 初始方案与最终方案\\\bottomrule
\end{tabularx}\end{table}
\begin{table}[H]\centering\small\caption{每组保存的方案与评价记录}\label{tab:b-artifacts}
\begin{tabularx}{\linewidth}{lYY}\toprule
记录 & 内容 & 用途\\\midrule
\texttt{final\_plan.json} & 操作映射、分核和同核列表 & 提交计划；问题三的固定B计划\\
\texttt{final\_selection.json} & 候选名、计划哈希、工期、搬运 & 核对配对双方的计划是否一致\\
候选\texttt{profile.json} & 核级展开、COPY及Spill统计 & 解释片上恢复和候选入围\\
候选\texttt{result.json} & 工期、流水线事件与搬运字段 & 同计划B/C配对及约束核验\\\bottomrule
\end{tabularx}\end{table}
结果文件对1500个组合键及其中2784条成功候选记录做了数值、计划和时间线一致性检查，另有15项跨场景抽样重算与存档结果逐项相同。

本文使用题面给出的100张计算图、硬件配置和评估程序。主表包含100图、1至5核和A、B、C三种场景的1500条记录；其中A的2至5核400条为正式问题一结果，另保存A单核100条用于核对整图A0基准（\texttt{baseline\_A0}），并单独标记单核搜索诊断（\texttt{search\_A\_K1}）。这两类记录用途不同：前者定义统一速度比基准，后者保留单核候选搜索的实际工期与比值。另有500组B/C同计划配对和24组Cache参数固定计划评价。表中工期、搬运和命中率均由相应评价结果汇总。

图中区间以100图有放回重采样1500次所得经验分位数表示实例组成对均值的影响。加速比均由原始工期计算。统一基准工期为整图单Task单核 A0 的 $T_{1,i}$；A、B 单核曲线按题面显示为 1.000，逐例表的 A/B 单核单元仍列各自最终已评计划的实际工期，C 单核单元列实际工期比。以\Case{053}为例，A/B/C 的单核实际工期分别为 1471202、1448004、1426552 cycle，统一基准为 1472686 cycle。

\section{三问逐图结果}\label{app:results}
表\ref{tab:detail-1}至表\ref{tab:detail-5}共列出100图、5种核数、3问的1500条最终记录。
每个$T_r/S_r$单元中的斜线分别分隔工期（cycle）和相对整图单核 A0 的加速比；$D_r$是官方额外搬运字节。A、B 的单核 $S_r$ 按题面显示为 1.000，C 的单核 $S_r$ 保留实际工期比；三种场景的单核实际工期均来自对应最终已评计划。
逐例表只列出最终工期、相对整图单核 A0 的加速比和额外数据搬运量。初始工期、候选来源、方案哈希、切分搬运与 Spill 搬运见 \path{data/audited_20260924/main_results.csv}；\path{data/final_idea_v2/candidate_effects.csv} 另列首个成功候选至最终候选的逐组工期变化。
\begingroup
\small
\setlength{\tabcolsep}{2.5pt}
\setlength{\LTleft}{\fill}\setlength{\LTright}{\fill}
\renewcommand{\arraystretch}{1.03}
\begin{longtable}{crrrrrr}
\caption{1核三问逐例工期/加速比与额外搬运（100图）}\label{tab:detail-1}\\
\toprule
图 & $T_A/S_A$ & $D_A$/B & $T_B/S_B$ & $D_B$/B & $T_C/S_C$ & $D_C$/B \\
\midrule
\endfirsthead
\multicolumn{7}{c}{续表}\\
\toprule
图 & $T_A/S_A$ & $D_A$/B & $T_B/S_B$ & $D_B$/B & $T_C/S_C$ & $D_C$/B \\
\midrule
\endhead
\midrule\multicolumn{7}{r}{续下页}\\
\endfoot
\bottomrule
\endlastfoot
001 & 233110/1.000 & 0 & 233110/1.000 & 1961984 & 233110/1.000 & 1961984 \\
002 & 261945/1.000 & 0 & 261945/1.000 & 645120 & 261945/1.000 & 0 \\
003 & 602212/1.000 & 0 & 592981/1.000 & 0 & 592981/1.016 & 0 \\
004 & 392357/1.000 & 0 & 392357/1.000 & 2856960 & 392357/1.000 & 663552 \\
005 & 95618/1.000 & 0 & 93910/1.000 & 0 & 93910/1.018 & 0 \\
006 & 82837/1.000 & 16400 & 76655/1.000 & 172068 & 76655/1.081 & 172068 \\
007 & 1012562/1.000 & 0 & 1012562/1.000 & 10152960 & 1012562/1.000 & 2810880 \\
008 & 487605/1.000 & 0 & 405761/1.000 & 7059072 & 335225/1.455 & 3520896 \\
009 & 153007/1.000 & 17606 & 153261/1.000 & 12288 & 153261/1.000 & 12288 \\
010 & 86969/1.000 & 0 & 82266/1.000 & 0 & 82266/1.057 & 0 \\
011 & 184697/1.000 & 0 & 184697/1.000 & 167936 & 184697/1.000 & 167936 \\
012 & 59088/1.000 & 0 & 50002/1.000 & 16384 & 49967/1.183 & 16384 \\
013 & 263949/1.000 & 0 & 263949/1.000 & 3872768 & 263949/1.000 & 933888 \\
014 & 17698327/1.000 & 71559072 & 17698626/1.000 & 278109264 & 17698626/1.000 & 278109264 \\
015 & 177362/1.000 & 180256 & 177362/1.000 & 502164 & 177362/1.000 & 180332 \\
016 & 7715523/1.000 & 0 & 7715523/1.000 & 0 & 7715523/1.000 & 0 \\
017 & 159480/1.000 & 14370 & 162810/1.000 & 315968 & 162810/1.000 & 315968 \\
018 & 1004367/1.000 & 1970176 & 1004367/1.000 & 6824960 & 1003110/1.001 & 1970176 \\
019 & 81144/1.000 & 0 & 70264/1.000 & 6144 & 70264/1.155 & 6144 \\
020 & 532746/1.000 & 0 & 532746/1.000 & 17174528 & 532746/1.000 & 12726272 \\
021 & 6520170/1.000 & 4702208 & 6567937/1.000 & 63094784 & 6567937/1.000 & 29097984 \\
022 & 56671/1.000 & 14094 & 54906/1.000 & 57344 & 54906/1.034 & 57344 \\
023 & 122614/1.000 & 0 & 116685/1.000 & 274948 & 116571/1.052 & 274948 \\
024 & 2555343/1.000 & 0 & 2555343/1.000 & 0 & 2555343/1.000 & 0 \\
025 & 4863026/1.000 & 0 & 4863026/1.000 & 56653824 & 4863026/1.000 & 17399808 \\
026 & 99203/1.000 & 0 & 95307/1.000 & 0 & 95307/1.041 & 0 \\
027 & 2325280/1.000 & 1425408 & 2341632/1.000 & 8091648 & 2341632/1.010 & 8091648 \\
028 & 43836040/1.000 & 277412256 & 43281233/1.000 & 201071064 & 42606375/1.029 & 201071064 \\
029 & 37555/1.000 & 0 & 37555/1.000 & 135216 & 37555/1.000 & 135216 \\
030 & 2807523/1.000 & 5197824 & 2808034/1.000 & 51834880 & 2808034/1.000 & 16946176 \\
031 & 1410637/1.000 & 3988480 & 1414369/1.000 & 16710144 & 1414369/1.000 & 16710144 \\
032 & 41443/1.000 & 0 & 34533/1.000 & 0 & 34533/1.200 & 0 \\
033 & 1773973/1.000 & 5191680 & 1774196/1.000 & 23242240 & 1774196/1.000 & 9474056 \\
034 & 838424/1.000 & 793104 & 842996/1.000 & 7961088 & 842733/1.000 & 1079832 \\
035 & 169571/1.000 & 0 & 169571/1.000 & 103472 & 169571/1.000 & 103472 \\
036 & 931510/1.000 & 0 & 931510/1.000 & 11882496 & 931510/1.000 & 6344704 \\
037 & 210001/1.000 & 0 & 210001/1.000 & 2600960 & 210001/1.000 & 2600960 \\
038 & 1363201/1.000 & 4233216 & 1364555/1.000 & 20147712 & 1364555/1.000 & 20147712 \\
039 & 1430004/1.000 & 7954432 & 1435262/1.000 & 28258304 & 1435262/1.000 & 28258304 \\
040 & 253454/1.000 & 0 & 253454/1.000 & 235008 & 253454/1.000 & 235008 \\
041 & 6003957/1.000 & 30898176 & 6003957/1.000 & 104178688 & 6003957/1.000 & 53859328 \\
042 & 1670133/1.000 & 6974464 & 1686740/1.000 & 17703424 & 1686740/1.000 & 17703424 \\
043 & 677656/1.000 & 49160 & 675739/1.000 & 78472 & 675739/1.003 & 78472 \\
044 & 94856/1.000 & 175616 & 92892/1.000 & 76288 & 92892/1.662 & 76288 \\
045 & 125614/1.000 & 0 & 125614/1.000 & 3690496 & 125614/1.000 & 2719744 \\
046 & 233229/1.000 & 1033216 & 227889/1.000 & 745472 & 227889/1.213 & 745472 \\
047 & 326508/1.000 & 0 & 324985/1.000 & 0 & 324985/1.005 & 0 \\
048 & 222220/1.000 & 0 & 222220/1.000 & 0 & 222220/1.000 & 0 \\
049 & 196276/1.000 & 0 & 196276/1.000 & 18432 & 196276/1.000 & 18432 \\
050 & 149690/1.000 & 0 & 149690/1.000 & 348192 & 149470/1.001 & 348192 \\
051 & 607628/1.000 & 0 & 607628/1.000 & 0 & 607628/1.000 & 0 \\
052 & 173557/1.000 & 16392 & 178673/1.000 & 225536 & 178673/1.009 & 225536 \\
053 & 1471202/1.000 & 5508790 & 1448004/1.000 & 4047952 & 1426552/1.032 & 4047952 \\
054 & 814055/1.000 & 81734 & 796718/1.000 & 165700 & 796718/1.023 & 165700 \\
055 & 116367/1.000 & 73732 & 118835/1.000 & 431696 & 118835/1.000 & 431696 \\
056 & 277285/1.000 & 0 & 277285/1.000 & 0 & 277285/1.000 & 0 \\
057 & 59366/1.000 & 77832 & 58953/1.000 & 319558 & 58106/1.028 & 319558 \\
058 & 4534566/1.000 & 29376512 & 4539717/1.000 & 127987712 & 4539717/1.000 & 127987712 \\
059 & 818974/1.000 & 2072576 & 768400/1.000 & 6178816 & 768400/1.079 & 6178816 \\
060 & 886258/1.000 & 2398344 & 888789/1.000 & 18983692 & 888789/1.000 & 18842440 \\
061 & 77191/1.000 & 0 & 77191/1.000 & 0 & 77191/1.000 & 0 \\
062 & 2854809/1.000 & 0 & 2854809/1.000 & 23152128 & 2854809/1.000 & 7315968 \\
063 & 1070397/1.000 & 0 & 1070397/1.000 & 7303680 & 1070397/1.000 & 1938432 \\
064 & 19116/1.000 & 0 & 19116/1.000 & 0 & 19116/1.000 & 0 \\
065 & 51959/1.000 & 8200 & 42821/1.000 & 16388 & 42821/1.239 & 16388 \\
066 & 358590/1.000 & 56438 & 358742/1.000 & 27648 & 358742/1.001 & 27648 \\
067 & 63834834/1.000 & 264449408 & 63782329/1.000 & 233595392 & 63620658/1.003 & 233595392 \\
068 & 360885/1.000 & 51910 & 360987/1.000 & 44480 & 360987/1.003 & 44480 \\
069 & 25111/1.000 & 68740 & 23559/1.000 & 0 & 23559/1.067 & 0 \\
070 & 116204/1.000 & 100352 & 117300/1.000 & 0 & 117300/1.000 & 0 \\
071 & 18919/1.000 & 0 & 18622/1.000 & 0 & 18337/1.032 & 0 \\
072 & 23897828/1.000 & 108613680 & 23897828/1.000 & 279466632 & 23897828/1.000 & 279466632 \\
073 & 10068454/1.000 & 131995904 & 10067403/1.000 & 133783808 & 10067403/1.001 & 133783808 \\
074 & 1679412/1.000 & 6336512 & 1682239/1.000 & 15523850 & 1682239/1.000 & 15523850 \\
075 & 1441809/1.000 & 0 & 1433562/1.000 & 0 & 1433562/1.006 & 0 \\
076 & 6176086/1.000 & 25961472 & 6185113/1.000 & 139435008 & 6185113/1.000 & 64754688 \\
077 & 495320/1.000 & 0 & 495320/1.000 & 2056384 & 495320/1.000 & 2056384 \\
078 & 2260779/1.000 & 10175328 & 2259940/1.000 & 10188024 & 2189192/1.049 & 10188024 \\
079 & 30952656/1.000 & 32555008 & 31052251/1.000 & 310804480 & 31052251/1.000 & 310804480 \\
080 & 292086/1.000 & 667648 & 292138/1.000 & 4014080 & 292138/1.000 & 4014080 \\
081 & 151795/1.000 & 0 & 151795/1.000 & 1704132 & 151795/1.000 & 1704132 \\
082 & 720361/1.000 & 170630 & 717008/1.000 & 119232 & 717008/1.006 & 119232 \\
083 & 1057993/1.000 & 13934688 & 1040648/1.000 & 12382336 & 930103/1.138 & 12382336 \\
084 & 2507105/1.000 & 9216 & 1762609/1.000 & 19562496 & 1552121/1.615 & 9877056 \\
085 & 3187049/1.000 & 227974 & 3187417/1.000 & 771328 & 3187417/1.000 & 771328 \\
086 & 118403/1.000 & 0 & 114843/1.000 & 22272 & 114843/1.031 & 22272 \\
087 & 3974071/1.000 & 7902720 & 3971407/1.000 & 8246464 & 3967115/1.002 & 8246464 \\
088 & 243623/1.000 & 0 & 240293/1.000 & 93124 & 240293/1.014 & 93124 \\
089 & 693092/1.000 & 0 & 693092/1.000 & 58368 & 693092/1.000 & 58368 \\
090 & 1092362/1.000 & 15562608 & 1090066/1.000 & 15781584 & 998222/1.180 & 15781584 \\
091 & 9186769/1.000 & 53947008 & 9186769/1.000 & 171968128 & 9186769/1.000 & 90629248 \\
092 & 4815554/1.000 & 66000096 & 4203197/1.000 & 19239168 & 4076456/1.181 & 19239168 \\
093 & 74205/1.000 & 0 & 74205/1.000 & 47104 & 74205/1.000 & 0 \\
094 & 92786/1.000 & 51200 & 92511/1.000 & 313344 & 92511/1.026 & 313344 \\
095 & 2084853/1.000 & 0 & 1651133/1.000 & 33294720 & 1316513/1.584 & 16706688 \\
096 & 83517/1.000 & 68352 & 80228/1.000 & 0 & 80228/1.049 & 0 \\
097 & 11491509/1.000 & 10584064 & 11385302/1.000 & 40259584 & 11365969/1.011 & 30957568 \\
098 & 177430/1.000 & 0 & 174782/1.000 & 385068 & 174782/1.015 & 385068 \\
099 & 1620966/1.000 & 5918720 & 1622999/1.000 & 15709568 & 1622999/1.000 & 15709568 \\
100 & 148339/1.000 & 5640 & 145067/1.000 & 0 & 145067/1.042 & 0 \\
\end{longtable}

\begin{longtable}{crrrrrr}
\caption{2核三问逐例工期/加速比与额外搬运（100图）}\label{tab:detail-2}\\
\toprule
图 & $T_A/S_A$ & $D_A$/B & $T_B/S_B$ & $D_B$/B & $T_C/S_C$ & $D_C$/B \\
\midrule
\endfirsthead
\multicolumn{7}{c}{续表}\\
\toprule
图 & $T_A/S_A$ & $D_A$/B & $T_B/S_B$ & $D_B$/B & $T_C/S_C$ & $D_C$/B \\
\midrule
\endhead
\midrule\multicolumn{7}{r}{续下页}\\
\endfoot
\bottomrule
\endlastfoot
001 & 116868/1.995 & 1152 & 116868/1.995 & 1152 & 116868/1.995 & 1152 \\
002 & 261945/1.000 & 0 & 208461/1.257 & 4680192 & 205287/1.276 & 4680192 \\
003 & 602212/1.000 & 0 & 562101/1.071 & 1950532 & 562101/1.071 & 1950532 \\
004 & 199074/1.971 & 110592 & 199074/1.971 & 110592 & 199074/1.971 & 110592 \\
005 & 95618/1.000 & 0 & 93910/1.018 & 0 & 93910/1.018 & 0 \\
006 & 42696/1.940 & 56200 & 42696/1.940 & 56200 & 42696/1.940 & 56200 \\
007 & 507494/1.995 & 27648 & 507494/1.995 & 27648 & 507494/1.995 & 27648 \\
008 & 244984/1.990 & 0 & 244266/1.996 & 0 & 244266/1.996 & 0 \\
009 & 92198/1.662 & 0 & 142414/1.076 & 367942 & 142414/1.076 & 367942 \\
010 & 46568/1.868 & 48384 & 46568/1.868 & 48384 & 46568/1.868 & 48384 \\
011 & 92646/1.994 & 180224 & 92646/1.994 & 180224 & 92646/1.994 & 180224 \\
012 & 30058/1.966 & 29376 & 30058/1.966 & 29376 & 30058/1.966 & 29376 \\
013 & 132198/1.997 & 8192 & 132198/1.997 & 8192 & 132198/1.997 & 8192 \\
014 & 8912981/1.986 & 75727800 & 8912981/1.986 & 75727800 & 8904766/1.988 & 75727800 \\
015 & 90509/1.960 & 199948 & 90509/1.960 & 199948 & 90509/1.960 & 199948 \\
016 & 7715523/1.000 & 0 & 7715523/1.000 & 0 & 7715523/1.000 & 0 \\
017 & 81772/1.991 & 87552 & 81772/1.991 & 87552 & 81772/1.991 & 87552 \\
018 & 504468/1.991 & 3253248 & 504468/1.991 & 3253248 & 502571/1.998 & 3253248 \\
019 & 41382/1.961 & 27168 & 41382/1.961 & 27168 & 41382/1.961 & 27168 \\
020 & 266580/1.998 & 0 & 266580/1.998 & 0 & 266580/1.998 & 0 \\
021 & 3194570/2.056 & 3174400 & 3194570/2.056 & 3174400 & 3193634/2.057 & 3174400 \\
022 & 28324/2.005 & 6912 & 28324/2.005 & 6912 & 28324/2.005 & 6912 \\
023 & 61501/1.994 & 84384 & 61501/1.994 & 84384 & 61501/1.994 & 84384 \\
024 & 2555343/1.000 & 0 & 2555343/1.000 & 0 & 2555343/1.000 & 0 \\
025 & 2431840/2.000 & 27648 & 2431840/2.000 & 27648 & 2431840/2.000 & 27648 \\
026 & 55057/1.802 & 21936 & 55057/1.802 & 21936 & 55057/1.802 & 21936 \\
027 & 1120962/2.109 & 1744896 & 1120962/2.109 & 1744896 & 1120962/2.109 & 1744896 \\
028 & 23638628/1.854 & 277473600 & 23638628/1.854 & 277473600 & 23615145/1.856 & 277473600 \\
029 & 19740/1.902 & 0 & 19740/1.902 & 0 & 19740/1.902 & 0 \\
030 & 1403594/2.001 & 7644160 & 1403594/2.001 & 7644160 & 1403449/2.001 & 7644160 \\
031 & 712596/1.985 & 5250048 & 712596/1.985 & 5250048 & 711944/1.987 & 5250048 \\
032 & 21040/1.970 & 6216 & 21040/1.970 & 6216 & 21040/1.970 & 6216 \\
033 & 889697/1.994 & 6664224 & 889697/1.994 & 6664224 & 887951/1.998 & 6664224 \\
034 & 421427/2.000 & 1984000 & 421427/2.000 & 1984000 & 421244/2.001 & 1984000 \\
035 & 169571/1.000 & 0 & 151830/1.117 & 732286 & 151830/1.117 & 732286 \\
036 & 466068/1.999 & 1152 & 466068/1.999 & 1152 & 466068/1.999 & 1152 \\
037 & 105374/1.993 & 204800 & 105374/1.993 & 204800 & 105374/1.993 & 204800 \\
038 & 674500/2.023 & 6790656 & 674500/2.023 & 6790656 & 674382/2.023 & 6790656 \\
039 & 708561/2.026 & 11067392 & 708561/2.026 & 11067392 & 705615/2.034 & 11067392 \\
040 & 137746/1.840 & 53760 & 137746/1.840 & 53760 & 137720/1.840 & 53760 \\
041 & 2990398/2.008 & 33669120 & 2990398/2.008 & 33669120 & 2988290/2.009 & 33669120 \\
042 & 849076/1.987 & 7413024 & 849076/1.987 & 7413024 & 845668/1.995 & 7413024 \\
043 & 353851/1.915 & 506880 & 353851/1.915 & 506880 & 353851/1.915 & 506880 \\
044 & 94856/1.628 & 175616 & 73924/2.089 & 22528 & 73924/2.089 & 22528 \\
045 & 63059/1.992 & 0 & 63059/1.992 & 0 & 63059/1.992 & 0 \\
046 & 153482/1.801 & 4017792 & 166578/1.660 & 3521632 & 137552/2.010 & 3521632 \\
047 & 326508/1.000 & 0 & 324985/1.005 & 0 & 324985/1.005 & 0 \\
048 & 222220/1.000 & 0 & 222220/1.000 & 0 & 222220/1.000 & 0 \\
049 & 196276/1.000 & 0 & 178903/1.097 & 729830 & 178903/1.097 & 729830 \\
050 & 149690/1.000 & 0 & 149158/1.004 & 731414 & 149158/1.004 & 731414 \\
051 & 607628/1.000 & 0 & 607628/1.000 & 0 & 607628/1.000 & 0 \\
052 & 90903/1.983 & 102912 & 90903/1.983 & 102912 & 90903/1.983 & 102912 \\
053 & 828335/1.778 & 5974464 & 828335/1.778 & 5974464 & 796880/1.848 & 5974464 \\
054 & 814055/1.001 & 81734 & 796718/1.023 & 165700 & 796718/1.023 & 165700 \\
055 & 61686/1.926 & 74240 & 61686/1.926 & 74240 & 61686/1.926 & 74240 \\
056 & 277285/1.000 & 0 & 277285/1.000 & 0 & 277285/1.000 & 0 \\
057 & 30183/1.979 & 22304 & 30183/1.979 & 22304 & 30183/1.979 & 22304 \\
058 & 2271316/1.999 & 33431552 & 2271316/1.999 & 33431552 & 2270206/2.000 & 33431552 \\
059 & 369318/2.246 & 360448 & 369318/2.246 & 360448 & 369318/2.246 & 360448 \\
060 & 443955/2.002 & 2603148 & 443955/2.002 & 2603148 & 443733/2.003 & 2603148 \\
061 & 42478/1.817 & 34560 & 42478/1.817 & 34560 & 42478/1.817 & 34560 \\
062 & 2854809/1.000 & 0 & 2610093/1.094 & 76270080 & 2525747/1.130 & 76270080 \\
063 & 795582/1.345 & 2116608 & 945557/1.132 & 26079744 & 914171/1.171 & 26079744 \\
064 & 19116/1.000 & 0 & 19116/1.000 & 0 & 19116/1.000 & 0 \\
065 & 26610/1.994 & 13968 & 26610/1.994 & 13968 & 26610/1.994 & 13968 \\
066 & 358590/1.001 & 56438 & 342680/1.048 & 1097258 & 342680/1.048 & 1097258 \\
067 & 63574293/1.004 & 224305664 & 33768704/1.890 & 264024576 & 33768704/1.890 & 264024576 \\
068 & 356477/1.016 & 1194322 & 339779/1.066 & 1310684 & 339318/1.067 & 1120508 \\
069 & 25111/1.001 & 68740 & 23559/1.067 & 0 & 23559/1.067 & 0 \\
070 & 60121/1.951 & 129536 & 60121/1.951 & 129536 & 60121/1.951 & 129536 \\
071 & 18919/1.000 & 0 & 18622/1.016 & 0 & 18337/1.032 & 0 \\
072 & 11872949/2.013 & 113068832 & 11872949/2.013 & 113068832 & 11862815/2.015 & 113068832 \\
073 & 9989226/1.009 & 114749952 & 5980423/1.686 & 141121024 & 5980423/1.686 & 141121024 \\
074 & 839999/2.003 & 6590472 & 839999/2.003 & 6590472 & 837941/2.008 & 6590472 \\
075 & 1441809/1.000 & 0 & 1433562/1.006 & 0 & 1433562/1.006 & 0 \\
076 & 3065214/2.018 & 31095808 & 3065214/2.018 & 31095808 & 3063843/2.019 & 31095808 \\
077 & 248734/1.991 & 57344 & 248734/1.991 & 57344 & 248734/1.991 & 57344 \\
078 & 1217774/1.886 & 12841128 & 1260358/1.822 & 12535344 & 1260358/1.822 & 12535344 \\
079 & 15446486/2.010 & 31780864 & 15446486/2.010 & 31780864 & 15446486/2.010 & 31780864 \\
080 & 145926/2.002 & 2191360 & 145926/2.002 & 2191360 & 145674/2.005 & 2191360 \\
081 & 76860/1.975 & 0 & 76860/1.975 & 0 & 76860/1.975 & 0 \\
082 & 718618/1.003 & 513296 & 717008/1.006 & 119232 & 717008/1.006 & 119232 \\
083 & 905273/1.169 & 2650240 & 775372/1.364 & 5429408 & 755398/1.401 & 5429408 \\
084 & 1255052/1.998 & 0 & 1255052/1.998 & 0 & 1255052/1.998 & 0 \\
085 & 3183732/1.001 & 719378 & 3187417/1.000 & 771328 & 3187417/1.000 & 771328 \\
086 & 118403/1.000 & 0 & 114843/1.031 & 22272 & 114843/1.031 & 22272 \\
087 & 2016479/1.971 & 10050048 & 2016479/1.971 & 10050048 & 2013301/1.974 & 10050048 \\
088 & 243623/1.000 & 0 & 240293/1.014 & 93124 & 240293/1.014 & 93124 \\
089 & 346832/1.998 & 238720 & 346832/1.998 & 238720 & 346832/1.998 & 238720 \\
090 & 1092362/1.079 & 15562608 & 836229/1.409 & 21941208 & 766154/1.538 & 21941208 \\
091 & 4571759/2.009 & 57904096 & 4571759/2.009 & 57904096 & 4569395/2.011 & 57904096 \\
092 & 2828252/1.703 & 66389088 & 2828252/1.703 & 66389088 & 2358868/2.041 & 66389088 \\
093 & 37827/1.962 & 8192 & 37827/1.962 & 8192 & 37827/1.962 & 8192 \\
094 & 48073/1.974 & 106272 & 48073/1.974 & 106272 & 48073/1.974 & 106272 \\
095 & 1042890/1.999 & 0 & 1042890/1.999 & 0 & 1042890/1.999 & 0 \\
096 & 44010/1.913 & 67392 & 44010/1.913 & 67392 & 44010/1.913 & 67392 \\
097 & 5617602/2.046 & 11796480 & 5617602/2.046 & 11796480 & 5617602/2.046 & 11796480 \\
098 & 90360/1.964 & 0 & 90360/1.964 & 0 & 90360/1.964 & 0 \\
099 & 808916/2.006 & 7254016 & 808916/2.006 & 7254016 & 806053/2.014 & 7254016 \\
100 & 81080/1.865 & 16896 & 81080/1.865 & 16896 & 81080/1.865 & 16896 \\
\end{longtable}

\begin{longtable}{crrrrrr}
\caption{3核三问逐例工期/加速比与额外搬运（100图）}\label{tab:detail-3}\\
\toprule
图 & $T_A/S_A$ & $D_A$/B & $T_B/S_B$ & $D_B$/B & $T_C/S_C$ & $D_C$/B \\
\midrule
\endfirsthead
\multicolumn{7}{c}{续表}\\
\toprule
图 & $T_A/S_A$ & $D_A$/B & $T_B/S_B$ & $D_B$/B & $T_C/S_C$ & $D_C$/B \\
\midrule
\endhead
\midrule\multicolumn{7}{r}{续下页}\\
\endfoot
\bottomrule
\endlastfoot
001 & 78545/2.968 & 2304 & 78545/2.968 & 2304 & 78545/2.968 & 2304 \\
002 & 173068/1.514 & 594432 & 208461/1.257 & 4680192 & 205287/1.276 & 4680192 \\
003 & 602212/1.000 & 0 & 562101/1.071 & 1950532 & 562101/1.071 & 1950532 \\
004 & 134043/2.927 & 221184 & 134043/2.927 & 221184 & 134043/2.927 & 221184 \\
005 & 95618/1.000 & 0 & 93910/1.018 & 0 & 93910/1.018 & 0 \\
006 & 31152/2.659 & 30080 & 31152/2.659 & 30080 & 31152/2.659 & 30080 \\
007 & 338094/2.995 & 55296 & 338094/2.995 & 55296 & 338094/2.995 & 55296 \\
008 & 164386/2.966 & 0 & 163359/2.985 & 0 & 163359/2.985 & 0 \\
009 & 92198/1.662 & 0 & 142414/1.076 & 367942 & 142414/1.076 & 367942 \\
010 & 32889/2.644 & 69120 & 32889/2.644 & 69120 & 32889/2.644 & 69120 \\
011 & 69625/2.653 & 1081344 & 69625/2.653 & 1081344 & 69625/2.653 & 1081344 \\
012 & 20661/2.860 & 39168 & 20661/2.860 & 39168 & 20661/2.860 & 39168 \\
013 & 88351/2.988 & 16384 & 88351/2.988 & 16384 & 88351/2.988 & 16384 \\
014 & 5993652/2.953 & 84088488 & 5993652/2.953 & 84088488 & 5957399/2.971 & 84088488 \\
015 & 62280/2.848 & 252416 & 62280/2.848 & 252416 & 62280/2.848 & 252416 \\
016 & 7715523/1.000 & 0 & 7715523/1.000 & 0 & 7715523/1.000 & 0 \\
017 & 55389/2.939 & 175104 & 55389/2.939 & 175104 & 55389/2.939 & 175104 \\
018 & 341047/2.945 & 4093952 & 341047/2.945 & 4093952 & 338488/2.967 & 4093952 \\
019 & 28657/2.832 & 54336 & 28657/2.832 & 54336 & 28657/2.832 & 54336 \\
020 & 177950/2.994 & 0 & 177950/2.994 & 0 & 177950/2.994 & 0 \\
021 & 2144024/3.063 & 7766016 & 2144024/3.063 & 7766016 & 2142851/3.065 & 7766016 \\
022 & 19026/2.984 & 13824 & 19026/2.984 & 13824 & 19026/2.984 & 13824 \\
023 & 41403/2.961 & 133536 & 41403/2.961 & 133536 & 41403/2.961 & 133536 \\
024 & 2555343/1.000 & 0 & 2555343/1.000 & 0 & 2555343/1.000 & 0 \\
025 & 1622164/2.998 & 55296 & 1622164/2.998 & 55296 & 1622164/2.998 & 55296 \\
026 & 41665/2.381 & 95664 & 41665/2.381 & 95664 & 41665/2.381 & 95664 \\
027 & 1120962/2.109 & 1744896 & 1120962/2.109 & 1744896 & 1120962/2.109 & 1744896 \\
028 & 16770711/2.614 & 277369056 & 16770711/2.614 & 277369056 & 16770711/2.614 & 277369056 \\
029 & 12743/2.947 & 0 & 12743/2.947 & 0 & 12743/2.947 & 0 \\
030 & 933857/3.007 & 9242624 & 933857/3.007 & 9242624 & 933236/3.009 & 9242624 \\
031 & 473822/2.985 & 7544832 & 473822/2.985 & 7544832 & 473000/2.990 & 7544832 \\
032 & 14175/2.924 & 12432 & 14175/2.924 & 12432 & 14175/2.924 & 12432 \\
033 & 592065/2.997 & 6990912 & 592065/2.997 & 6990912 & 590152/3.006 & 6990912 \\
034 & 276994/3.043 & 2830336 & 276994/3.043 & 2830336 & 276560/3.048 & 2830336 \\
035 & 169571/1.000 & 0 & 151830/1.117 & 732286 & 151830/1.117 & 732286 \\
036 & 311345/2.992 & 2304 & 311345/2.992 & 2304 & 311345/2.992 & 2304 \\
037 & 74038/2.836 & 1064960 & 74038/2.836 & 1064960 & 74038/2.836 & 1064960 \\
038 & 452248/3.017 & 8097792 & 452248/3.017 & 8097792 & 452022/3.019 & 8097792 \\
039 & 485565/2.956 & 14323712 & 485565/2.956 & 14323712 & 483226/2.970 & 14323712 \\
040 & 110105/2.302 & 0 & 110105/2.302 & 0 & 110105/2.302 & 0 \\
041 & 1991911/3.014 & 34851840 & 1991911/3.014 & 34851840 & 1989953/3.017 & 34851840 \\
042 & 562868/2.997 & 8334912 & 562868/2.997 & 8334912 & 561628/3.003 & 8334912 \\
043 & 268777/2.521 & 141320 & 268777/2.521 & 141320 & 268777/2.521 & 141320 \\
044 & 94856/1.628 & 175616 & 73924/2.089 & 22528 & 73924/2.089 & 22528 \\
045 & 42601/2.949 & 0 & 42601/2.949 & 0 & 42601/2.949 & 0 \\
046 & 142176/1.944 & 4021280 & 142176/1.944 & 4021280 & 121494/2.275 & 4021280 \\
047 & 326508/1.000 & 0 & 324985/1.005 & 0 & 324985/1.005 & 0 \\
048 & 222220/1.000 & 0 & 222220/1.000 & 0 & 222220/1.000 & 0 \\
049 & 196276/1.000 & 0 & 178903/1.097 & 729830 & 178903/1.097 & 729830 \\
050 & 149690/1.000 & 0 & 149158/1.004 & 731414 & 149158/1.004 & 731414 \\
051 & 607628/1.000 & 0 & 607628/1.000 & 0 & 607628/1.000 & 0 \\
052 & 61071/2.952 & 181248 & 61071/2.952 & 181248 & 61071/2.952 & 181248 \\
053 & 542159/2.716 & 5803072 & 828335/1.778 & 5974464 & 796880/1.848 & 5974464 \\
054 & 814055/1.001 & 81734 & 796718/1.023 & 165700 & 796718/1.023 & 165700 \\
055 & 42047/2.826 & 115712 & 42047/2.826 & 115712 & 42047/2.826 & 115712 \\
056 & 277285/1.000 & 0 & 277285/1.000 & 0 & 277285/1.000 & 0 \\
057 & 20705/2.885 & 14112 & 20705/2.885 & 14112 & 20705/2.885 & 14112 \\
058 & 1554831/2.920 & 30400512 & 1554831/2.920 & 30400512 & 1554831/2.920 & 30400512 \\
059 & 266217/3.116 & 2998272 & 266217/3.116 & 2998272 & 262333/3.162 & 2998272 \\
060 & 298210/2.980 & 3221760 & 298210/2.980 & 3221760 & 297398/2.989 & 3221760 \\
061 & 30303/2.547 & 36096 & 30303/2.547 & 36096 & 30303/2.547 & 36096 \\
062 & 2854809/1.000 & 0 & 2610093/1.094 & 76270080 & 2525747/1.130 & 76270080 \\
063 & 795582/1.345 & 2116608 & 945557/1.132 & 26079744 & 914171/1.171 & 26079744 \\
064 & 19116/1.000 & 0 & 19116/1.000 & 0 & 19116/1.000 & 0 \\
065 & 18032/2.943 & 27936 & 18032/2.943 & 27936 & 18032/2.943 & 27936 \\
066 & 358549/1.001 & 429368 & 316559/1.134 & 727228 & 316556/1.134 & 727228 \\
067 & 23568859/2.708 & 263599744 & 23568859/2.708 & 263599744 & 23568859/2.708 & 263599744 \\
068 & 356477/1.016 & 1194322 & 339779/1.066 & 1310684 & 339318/1.067 & 1120508 \\
069 & 25111/1.001 & 68740 & 23559/1.067 & 0 & 23559/1.067 & 0 \\
070 & 39333/2.982 & 259072 & 39333/2.982 & 259072 & 39333/2.982 & 259072 \\
071 & 18919/1.000 & 0 & 18622/1.016 & 0 & 18337/1.032 & 0 \\
072 & 7935290/3.012 & 120023840 & 7935290/3.012 & 120023840 & 7922729/3.016 & 120023840 \\
073 & 4515227/2.233 & 140671616 & 4515227/2.233 & 140671616 & 4515227/2.233 & 140671616 \\
074 & 563098/2.987 & 7639040 & 563098/2.987 & 7639040 & 557771/3.016 & 7639040 \\
075 & 1441809/1.000 & 0 & 1433562/1.006 & 0 & 1433562/1.006 & 0 \\
076 & 2050028/3.017 & 34332672 & 2050028/3.017 & 34332672 & 2047868/3.020 & 34332672 \\
077 & 168976/2.931 & 114688 & 168976/2.931 & 114688 & 168976/2.931 & 114688 \\
078 & 1217774/1.886 & 12841128 & 1260358/1.822 & 12535344 & 1260358/1.822 & 12535344 \\
079 & 10213492/3.040 & 32980992 & 10181692/3.050 & 39960576 & 10175315/3.052 & 39960576 \\
080 & 101563/2.876 & 1531904 & 100325/2.912 & 3047424 & 100325/2.912 & 3047424 \\
081 & 51943/2.922 & 0 & 51943/2.922 & 0 & 51943/2.922 & 0 \\
082 & 718618/1.003 & 513296 & 691651/1.042 & 1474314 & 691651/1.042 & 1474314 \\
083 & 512623/2.064 & 14842752 & 512623/2.064 & 14842752 & 402691/2.627 & 14842752 \\
084 & 836220/2.999 & 0 & 836220/2.999 & 0 & 836220/2.999 & 0 \\
085 & 3183732/1.001 & 719378 & 3130505/1.018 & 2382092 & 3130505/1.018 & 2382092 \\
086 & 118403/1.000 & 0 & 114843/1.031 & 22272 & 114843/1.031 & 22272 \\
087 & 1374901/2.890 & 11036160 & 1374901/2.890 & 11036160 & 1372451/2.896 & 11036160 \\
088 & 243623/1.000 & 0 & 240293/1.014 & 93124 & 240293/1.014 & 93124 \\
089 & 243614/2.845 & 477440 & 243614/2.845 & 477440 & 243614/2.845 & 477440 \\
090 & 1003439/1.174 & 5735760 & 761694/1.547 & 20499936 & 706360/1.668 & 20499936 \\
091 & 3090094/2.973 & 64602944 & 3090094/2.973 & 64602944 & 3070819/2.992 & 64602944 \\
092 & 2231804/2.158 & 66784736 & 2231804/2.158 & 66784736 & 1792534/2.686 & 66784736 \\
093 & 25103/2.956 & 16384 & 25103/2.956 & 16384 & 25103/2.956 & 16384 \\
094 & 32222/2.945 & 212544 & 32222/2.945 & 212544 & 32222/2.945 & 212544 \\
095 & 696802/2.992 & 0 & 695775/2.996 & 0 & 695775/2.996 & 0 \\
096 & 30318/2.777 & 79488 & 30318/2.777 & 79488 & 30318/2.777 & 79488 \\
097 & 5617602/2.046 & 11796480 & 5617602/2.046 & 11796480 & 5617602/2.046 & 11796480 \\
098 & 89825/1.975 & 0 & 89825/1.975 & 0 & 89825/1.975 & 0 \\
099 & 544403/2.981 & 8798208 & 544403/2.981 & 8798208 & 540586/3.002 & 8798208 \\
100 & 57811/2.616 & 72192 & 57811/2.616 & 72192 & 57811/2.616 & 72192 \\
\end{longtable}

\begin{longtable}{crrrrrr}
\caption{4核三问逐例工期/加速比与额外搬运（100图）}\label{tab:detail-4}\\
\toprule
图 & $T_A/S_A$ & $D_A$/B & $T_B/S_B$ & $D_B$/B & $T_C/S_C$ & $D_C$/B \\
\midrule
\endfirsthead
\multicolumn{7}{c}{续表}\\
\toprule
图 & $T_A/S_A$ & $D_A$/B & $T_B/S_B$ & $D_B$/B & $T_C/S_C$ & $D_C$/B \\
\midrule
\endhead
\midrule\multicolumn{7}{r}{续下页}\\
\endfoot
\bottomrule
\endlastfoot
001 & 58984/3.952 & 3456 & 58984/3.952 & 3456 & 58984/3.952 & 3456 \\
002 & 173068/1.514 & 594432 & 129903/2.016 & 5766144 & 124800/2.099 & 5766144 \\
003 & 600785/1.002 & 829732 & 562101/1.071 & 1950532 & 562101/1.071 & 1950532 \\
004 & 134043/2.927 & 221184 & 104188/3.766 & 331776 & 104188/3.766 & 331776 \\
005 & 95618/1.000 & 0 & 93910/1.018 & 0 & 93910/1.018 & 0 \\
006 & 31152/2.659 & 30080 & 31152/2.659 & 30080 & 31152/2.659 & 30080 \\
007 & 254398/3.980 & 82944 & 254398/3.980 & 82944 & 254398/3.980 & 82944 \\
008 & 123060/3.962 & 0 & 123060/3.962 & 0 & 123060/3.962 & 0 \\
009 & 92198/1.662 & 0 & 142414/1.076 & 367942 & 142414/1.076 & 367942 \\
010 & 29918/2.907 & 96768 & 29918/2.907 & 96768 & 29858/2.913 & 96768 \\
011 & 46880/3.940 & 540672 & 46880/3.940 & 540672 & 46880/3.940 & 540672 \\
012 & 15863/3.725 & 49536 & 15863/3.725 & 49536 & 15863/3.725 & 49536 \\
013 & 66982/3.941 & 24576 & 66982/3.941 & 24576 & 66982/3.941 & 24576 \\
014 & 4551951/3.888 & 90028440 & 4551951/3.888 & 90028440 & 4526148/3.910 & 90028440 \\
015 & 48727/3.640 & 362248 & 48727/3.640 & 362248 & 48669/3.644 & 362248 \\
016 & 7715523/1.000 & 0 & 7715523/1.000 & 0 & 7715523/1.000 & 0 \\
017 & 41209/3.951 & 253440 & 41209/3.951 & 253440 & 41209/3.951 & 253440 \\
018 & 253676/3.959 & 4938752 & 253676/3.959 & 4938752 & 253384/3.964 & 4938752 \\
019 & 24099/3.367 & 54336 & 24099/3.367 & 54336 & 24099/3.367 & 54336 \\
020 & 133704/3.985 & 0 & 133704/3.985 & 0 & 133704/3.985 & 0 \\
021 & 1550324/4.236 & 4608000 & 1550324/4.236 & 4608000 & 1550324/4.236 & 4608000 \\
022 & 14526/3.909 & 13824 & 14526/3.909 & 13824 & 14526/3.909 & 13824 \\
023 & 30142/4.068 & 281280 & 30142/4.068 & 281280 & 30142/4.068 & 281280 \\
024 & 2555343/1.000 & 0 & 2555343/1.000 & 0 & 2555343/1.000 & 0 \\
025 & 1217444/3.994 & 82944 & 1217444/3.994 & 82944 & 1217444/3.994 & 82944 \\
026 & 41665/2.381 & 95664 & 41665/2.381 & 95664 & 41665/2.381 & 95664 \\
027 & 560838/4.215 & 2949120 & 559196/4.227 & 3317760 & 559196/4.227 & 3317760 \\
028 & 13608312/3.221 & 277760640 & 13608312/3.221 & 277760640 & 13585418/3.227 & 277760640 \\
029 & 11028/3.405 & 0 & 11028/3.405 & 0 & 11028/3.405 & 0 \\
030 & 700342/4.010 & 10404864 & 700342/4.010 & 10404864 & 700156/4.011 & 10404864 \\
031 & 357987/3.951 & 7953408 & 357987/3.951 & 7953408 & 355839/3.975 & 7953408 \\
032 & 11004/3.766 & 18576 & 11004/3.766 & 18576 & 11004/3.766 & 18576 \\
033 & 443447/4.001 & 7214688 & 443447/4.001 & 7214688 & 442667/4.008 & 7214688 \\
034 & 212170/3.973 & 2978304 & 212170/3.973 & 2978304 & 212170/3.973 & 2978304 \\
035 & 169571/1.000 & 0 & 151830/1.117 & 732286 & 151830/1.117 & 732286 \\
036 & 233584/3.988 & 3456 & 233584/3.988 & 3456 & 233584/3.988 & 3456 \\
037 & 53320/3.939 & 614400 & 53320/3.939 & 614400 & 53320/3.939 & 614400 \\
038 & 339624/4.018 & 9295872 & 339624/4.018 & 9295872 & 337813/4.039 & 9295872 \\
039 & 395323/3.631 & 17604608 & 395323/3.631 & 17604608 & 368055/3.900 & 17604608 \\
040 & 110105/2.302 & 0 & 110105/2.302 & 0 & 110105/2.302 & 0 \\
041 & 1499144/4.005 & 39094272 & 1499144/4.005 & 39094272 & 1494347/4.018 & 39094272 \\
042 & 426511/3.955 & 8553312 & 426511/3.955 & 8553312 & 423652/3.981 & 8553312 \\
043 & 268777/2.521 & 141320 & 268777/2.521 & 141320 & 268777/2.521 & 141320 \\
044 & 94856/1.628 & 175616 & 73924/2.089 & 22528 & 73924/2.089 & 22528 \\
045 & 31878/3.940 & 0 & 31878/3.940 & 0 & 31878/3.940 & 0 \\
046 & 142176/1.944 & 4021280 & 142176/1.944 & 4021280 & 121494/2.275 & 4021280 \\
047 & 326508/1.000 & 0 & 324985/1.005 & 0 & 324985/1.005 & 0 \\
048 & 222220/1.000 & 0 & 222220/1.000 & 0 & 222220/1.000 & 0 \\
049 & 196276/1.000 & 0 & 178903/1.097 & 729830 & 178903/1.097 & 729830 \\
050 & 149690/1.000 & 0 & 149158/1.004 & 731414 & 149158/1.004 & 731414 \\
051 & 607628/1.000 & 0 & 607628/1.000 & 0 & 607628/1.000 & 0 \\
052 & 45610/3.953 & 254976 & 45610/3.953 & 254976 & 45610/3.953 & 254976 \\
053 & 492232/2.992 & 6360992 & 828335/1.778 & 5974464 & 796880/1.848 & 5974464 \\
054 & 814055/1.001 & 81734 & 796718/1.023 & 165700 & 796718/1.023 & 165700 \\
055 & 32099/3.702 & 173568 & 32099/3.702 & 173568 & 32099/3.702 & 173568 \\
056 & 277285/1.000 & 0 & 277285/1.000 & 0 & 277285/1.000 & 0 \\
057 & 16485/3.623 & 28224 & 16485/3.623 & 28224 & 16485/3.623 & 28224 \\
058 & 1154519/3.932 & 42622976 & 1154519/3.932 & 42622976 & 1142629/3.973 & 42622976 \\
059 & 185216/4.478 & 1081344 & 185216/4.478 & 1081344 & 185216/4.478 & 1081344 \\
060 & 221720/4.009 & 3930376 & 221720/4.009 & 3930376 & 221694/4.009 & 3930376 \\
061 & 24024/3.213 & 70656 & 24024/3.213 & 70656 & 24024/3.213 & 70656 \\
062 & 2537194/1.125 & 6633984 & 2240290/1.274 & 85255680 & 2195265/1.300 & 85255680 \\
063 & 348902/3.068 & 2469888 & 742037/1.443 & 28698624 & 742037/1.443 & 28698624 \\
064 & 19116/1.000 & 0 & 19116/1.000 & 0 & 19116/1.000 & 0 \\
065 & 13541/3.919 & 41904 & 13541/3.919 & 41904 & 13541/3.919 & 41904 \\
066 & 358549/1.001 & 429368 & 316559/1.134 & 727228 & 316556/1.134 & 727228 \\
067 & 18473155/3.456 & 263174912 & 18473155/3.456 & 263174912 & 18473155/3.456 & 263174912 \\
068 & 356477/1.016 & 1194322 & 339779/1.066 & 1310684 & 339318/1.067 & 1120508 \\
069 & 25111/1.001 & 68740 & 23559/1.067 & 0 & 23559/1.067 & 0 \\
070 & 31225/3.757 & 360960 & 31225/3.757 & 360960 & 31192/3.761 & 360960 \\
071 & 18919/1.000 & 0 & 18622/1.016 & 0 & 18337/1.032 & 0 \\
072 & 6003776/3.980 & 123642544 & 6003776/3.980 & 123642544 & 5991003/3.989 & 123642544 \\
073 & 3903787/2.583 & 140222208 & 3903787/2.583 & 140222208 & 3903787/2.583 & 140222208 \\
074 & 418987/4.015 & 7852032 & 418987/4.015 & 7852032 & 415633/4.047 & 7852032 \\
075 & 1441809/1.000 & 0 & 1430810/1.008 & 1350930 & 1430808/1.008 & 1350930 \\
076 & 1540201/4.016 & 37442560 & 1540201/4.016 & 37442560 & 1536306/4.026 & 37442560 \\
077 & 127832/3.875 & 114688 & 127832/3.875 & 114688 & 127591/3.882 & 114688 \\
078 & 637384/3.602 & 14566248 & 690924/3.323 & 13146912 & 690924/3.323 & 13146912 \\
079 & 7635624/4.067 & 30445568 & 7635624/4.067 & 30445568 & 7635624/4.067 & 30445568 \\
080 & 101563/2.876 & 1531904 & 100325/2.912 & 3047424 & 90584/3.225 & 4190208 \\
081 & 39588/3.834 & 0 & 39588/3.834 & 0 & 39588/3.834 & 0 \\
082 & 718618/1.003 & 513296 & 691651/1.042 & 1474314 & 691651/1.042 & 1474314 \\
083 & 450029/2.351 & 14977600 & 450029/2.351 & 14977600 & 333151/3.176 & 14977600 \\
084 & 629028/3.986 & 0 & 629028/3.986 & 0 & 629028/3.986 & 0 \\
085 & 3178812/1.003 & 1688746 & 3130505/1.018 & 2382092 & 3130505/1.018 & 2382092 \\
086 & 118403/1.000 & 0 & 114843/1.031 & 22272 & 114843/1.031 & 22272 \\
087 & 1046846/3.796 & 12607488 & 1046846/3.796 & 12607488 & 1043492/3.808 & 12607488 \\
088 & 243623/1.000 & 0 & 240293/1.014 & 93124 & 240293/1.014 & 93124 \\
089 & 178160/3.890 & 716160 & 178160/3.890 & 716160 & 178160/3.890 & 716160 \\
090 & 608275/1.937 & 21392136 & 731084/1.612 & 19028400 & 674035/1.748 & 19028400 \\
091 & 2315897/3.967 & 69360288 & 2315897/3.967 & 69360288 & 2307637/3.981 & 69360288 \\
092 & 1927728/2.498 & 67180384 & 1927728/2.498 & 67180384 & 1470508/3.275 & 67180384 \\
093 & 19779/3.752 & 24576 & 19779/3.752 & 24576 & 19779/3.752 & 24576 \\
094 & 24640/3.851 & 318528 & 24640/3.851 & 318528 & 24640/3.851 & 318528 \\
095 & 524422/3.976 & 0 & 524422/3.976 & 0 & 524422/3.976 & 0 \\
096 & 26984/3.120 & 79488 & 26984/3.120 & 79488 & 26984/3.120 & 79488 \\
097 & 2626908/4.375 & 14221312 & 2626908/4.375 & 14221312 & 2626908/4.375 & 14221312 \\
098 & 89825/1.975 & 0 & 89825/1.975 & 0 & 89825/1.975 & 0 \\
099 & 405336/4.004 & 9838592 & 405336/4.004 & 9838592 & 403555/4.022 & 9838592 \\
100 & 57811/2.616 & 72192 & 57811/2.616 & 72192 & 57811/2.616 & 72192 \\
\end{longtable}

\begin{longtable}{crrrrrr}
\caption{5核三问逐例工期/加速比与额外搬运（100图）}\label{tab:detail-5}\\
\toprule
图 & $T_A/S_A$ & $D_A$/B & $T_B/S_B$ & $D_B$/B & $T_C/S_C$ & $D_C$/B \\
\midrule
\endfirsthead
\multicolumn{7}{c}{续表}\\
\toprule
图 & $T_A/S_A$ & $D_A$/B & $T_B/S_B$ & $D_B$/B & $T_C/S_C$ & $D_C$/B \\
\midrule
\endhead
\midrule\multicolumn{7}{r}{续下页}\\
\endfoot
\bottomrule
\endlastfoot
001 & 47502/4.907 & 4608 & 47502/4.907 & 4608 & 47502/4.907 & 4608 \\
002 & 128412/2.040 & 700416 & 129903/2.016 & 5766144 & 124800/2.099 & 5766144 \\
003 & 600785/1.002 & 829732 & 541371/1.112 & 2325078 & 541371/1.112 & 2325078 \\
004 & 88341/4.441 & 442368 & 88341/4.441 & 442368 & 88341/4.441 & 442368 \\
005 & 95618/1.000 & 0 & 93910/1.018 & 0 & 93910/1.018 & 0 \\
006 & 31152/2.659 & 30080 & 31152/2.659 & 30080 & 31152/2.659 & 30080 \\
007 & 203834/4.968 & 110592 & 203834/4.968 & 110592 & 203834/4.968 & 110592 \\
008 & 100603/4.847 & 0 & 100603/4.847 & 0 & 100603/4.847 & 0 \\
009 & 92198/1.662 & 0 & 139951/1.095 & 1098596 & 138757/1.105 & 1098596 \\
010 & 29918/2.907 & 96768 & 29918/2.907 & 96768 & 29858/2.913 & 96768 \\
011 & 46880/3.940 & 540672 & 46880/3.940 & 540672 & 46777/3.948 & 1622016 \\
012 & 13406/4.408 & 53568 & 13406/4.408 & 53568 & 13406/4.408 & 53568 \\
013 & 54363/4.855 & 32768 & 54363/4.855 & 32768 & 54363/4.855 & 32768 \\
014 & 3712893/4.767 & 93975048 & 3712893/4.767 & 93975048 & 3700752/4.782 & 93975048 \\
015 & 43665/4.062 & 336384 & 43665/4.062 & 336384 & 42872/4.137 & 336384 \\
016 & 7715523/1.000 & 0 & 7715523/1.000 & 0 & 7715523/1.000 & 0 \\
017 & 34231/4.756 & 338688 & 34231/4.756 & 338688 & 34231/4.756 & 338688 \\
018 & 205219/4.894 & 5345280 & 205219/4.894 & 5345280 & 201659/4.981 & 5345280 \\
019 & 24099/3.367 & 54336 & 24099/3.367 & 54336 & 24099/3.367 & 54336 \\
020 & 107650/4.949 & 0 & 107650/4.949 & 0 & 107650/4.949 & 0 \\
021 & 1243902/5.280 & 7372800 & 1250168/5.254 & 11431936 & 1249380/5.257 & 11431936 \\
022 & 11824/4.802 & 20736 & 11824/4.802 & 20736 & 11824/4.802 & 20736 \\
023 & 24517/5.001 & 270912 & 24517/5.001 & 270912 & 24517/5.001 & 270912 \\
024 & 2555343/1.000 & 0 & 2555343/1.000 & 0 & 2555343/1.000 & 0 \\
025 & 975282/4.986 & 110592 & 975282/4.986 & 110592 & 975282/4.986 & 110592 \\
026 & 41665/2.381 & 95664 & 41665/2.381 & 95664 & 41655/2.382 & 21936 \\
027 & 558753/4.231 & 5529600 & 558753/4.231 & 5529600 & 558753/4.231 & 5529600 \\
028 & 11491035/3.815 & 277818912 & 11491035/3.815 & 277818912 & 11450794/3.828 & 277818912 \\
029 & 9383/4.002 & 0 & 9383/4.002 & 0 & 9383/4.002 & 0 \\
030 & 561102/5.004 & 11479040 & 561102/5.004 & 11479040 & 560897/5.006 & 11479040 \\
031 & 285693/4.951 & 7931904 & 285693/4.951 & 7931904 & 284489/4.972 & 7931904 \\
032 & 10496/3.948 & 12432 & 10496/3.948 & 12432 & 10479/3.955 & 12432 \\
033 & 355102/4.996 & 7816320 & 355102/4.996 & 7816320 & 354514/5.005 & 7816320 \\
034 & 170371/4.948 & 3621888 & 170371/4.948 & 3621888 & 168925/4.990 & 3621888 \\
035 & 95460/1.776 & 4608 & 151830/1.117 & 732286 & 151830/1.117 & 732286 \\
036 & 187182/4.976 & 4608 & 187182/4.976 & 4608 & 187182/4.976 & 4608 \\
037 & 44206/4.751 & 1515520 & 43800/4.795 & 1925120 & 43013/4.882 & 1925120 \\
038 & 275133/4.960 & 9676800 & 275133/4.960 & 9676800 & 272533/5.007 & 9676800 \\
039 & 315610/4.548 & 9457664 & 395323/3.631 & 17604608 & 340476/4.215 & 20480000 \\
040 & 110105/2.302 & 0 & 110105/2.302 & 0 & 110105/2.302 & 0 \\
041 & 1201884/4.995 & 40673280 & 1201884/4.995 & 40673280 & 1194643/5.026 & 40673280 \\
042 & 342520/4.925 & 10187328 & 342520/4.925 & 10187328 & 339457/4.969 & 10187328 \\
043 & 268777/2.521 & 141320 & 268777/2.521 & 141320 & 268777/2.521 & 141320 \\
044 & 94856/1.628 & 175616 & 73924/2.089 & 22528 & 73924/2.089 & 22528 \\
045 & 26164/4.801 & 0 & 26164/4.801 & 0 & 26164/4.801 & 0 \\
046 & 142176/1.944 & 4021280 & 142176/1.944 & 4021280 & 121494/2.275 & 4021280 \\
047 & 326508/1.000 & 0 & 316261/1.032 & 141400 & 316261/1.032 & 141400 \\
048 & 222220/1.000 & 0 & 222220/1.000 & 0 & 222220/1.000 & 0 \\
049 & 196276/1.000 & 0 & 178903/1.097 & 729830 & 178903/1.097 & 729830 \\
050 & 149690/1.000 & 0 & 147954/1.012 & 1893444 & 147851/1.012 & 1893444 \\
051 & 607628/1.000 & 0 & 607628/1.000 & 0 & 607628/1.000 & 0 \\
052 & 37387/4.823 & 328704 & 37387/4.823 & 328704 & 37387/4.823 & 328704 \\
053 & 492232/2.992 & 6360992 & 828335/1.778 & 5974464 & 796880/1.848 & 5974464 \\
054 & 814055/1.001 & 81734 & 796718/1.023 & 165700 & 796718/1.023 & 165700 \\
055 & 32099/3.702 & 173568 & 25635/4.636 & 231424 & 25635/4.636 & 231424 \\
056 & 277285/1.000 & 0 & 277285/1.000 & 0 & 277285/1.000 & 0 \\
057 & 14028/4.258 & 14112 & 14028/4.258 & 14112 & 14028/4.258 & 14112 \\
058 & 1004819/4.518 & 32104448 & 1004819/4.518 & 32104448 & 1004819/4.518 & 32104448 \\
059 & 150782/5.501 & 4325376 & 150782/5.501 & 4325376 & 150736/5.503 & 4325376 \\
060 & 175702/5.059 & 3948928 & 175702/5.059 & 3948928 & 175626/5.061 & 3948928 \\
061 & 22283/3.464 & 82944 & 22283/3.464 & 82944 & 22283/3.464 & 82944 \\
062 & 1661193/1.719 & 7405056 & 1723314/1.657 & 80862720 & 1690862/1.688 & 80862720 \\
063 & 269453/3.972 & 2755584 & 612094/1.749 & 29111808 & 602954/1.775 & 29111808 \\
064 & 19116/1.000 & 0 & 19116/1.000 & 0 & 19116/1.000 & 0 \\
065 & 11396/4.657 & 55872 & 11396/4.657 & 55872 & 11396/4.657 & 55872 \\
066 & 358549/1.001 & 429368 & 300629/1.194 & 2679682 & 300603/1.194 & 2679682 \\
067 & 18473155/3.456 & 263174912 & 15931687/4.007 & 262750080 & 15931687/4.007 & 262750080 \\
068 & 356477/1.016 & 1194322 & 339779/1.066 & 1310684 & 339318/1.067 & 1120508 \\
069 & 25111/1.001 & 68740 & 23559/1.067 & 0 & 23559/1.067 & 0 \\
070 & 26275/4.464 & 449024 & 26275/4.464 & 449024 & 26216/4.474 & 449024 \\
071 & 18919/1.000 & 0 & 18622/1.016 & 0 & 18337/1.032 & 0 \\
072 & 4831281/4.946 & 130884272 & 4831281/4.946 & 130884272 & 4809827/4.969 & 130884272 \\
073 & 3517515/2.866 & 139772800 & 3517515/2.866 & 139772800 & 3517515/2.866 & 139772800 \\
074 & 337275/4.988 & 7983104 & 337275/4.988 & 7983104 & 335626/5.012 & 7983104 \\
075 & 1441809/1.000 & 0 & 1430487/1.008 & 430232 & 1430487/1.008 & 430232 \\
076 & 1240011/4.988 & 40553472 & 1240011/4.988 & 40553472 & 1232435/5.019 & 40553472 \\
077 & 115617/4.284 & 172032 & 115617/4.284 & 172032 & 115617/4.284 & 172032 \\
078 & 637384/3.602 & 14566248 & 690924/3.323 & 13146912 & 690924/3.323 & 13146912 \\
079 & 6075976/5.111 & 47677440 & 6075976/5.111 & 47677440 & 6075472/5.111 & 47677440 \\
080 & 101563/2.876 & 1531904 & 100325/2.912 & 3047424 & 90408/3.231 & 4571136 \\
081 & 30943/4.906 & 0 & 30943/4.906 & 0 & 30943/4.906 & 0 \\
082 & 718618/1.003 & 513296 & 691651/1.042 & 1474314 & 691651/1.042 & 1474314 \\
083 & 442480/2.391 & 15058048 & 450029/2.351 & 14977600 & 333151/3.176 & 14977600 \\
084 & 503472/4.980 & 0 & 503472/4.980 & 0 & 503472/4.980 & 0 \\
085 & 3170837/1.005 & 2118612 & 3130505/1.018 & 2382092 & 3130505/1.018 & 2382092 \\
086 & 118403/1.000 & 0 & 114843/1.031 & 22272 & 114843/1.031 & 22272 \\
087 & 833349/4.769 & 12487680 & 833349/4.769 & 12487680 & 829669/4.790 & 12487680 \\
088 & 243623/1.000 & 0 & 240293/1.014 & 93124 & 240293/1.014 & 93124 \\
089 & 147432/4.701 & 954880 & 147432/4.701 & 954880 & 147432/4.701 & 954880 \\
090 & 542638/2.172 & 21724704 & 542638/2.172 & 21724704 & 509924/2.311 & 21724704 \\
091 & 1882479/4.880 & 73417216 & 1882479/4.880 & 73417216 & 1876653/4.895 & 73417216 \\
092 & 1661234/2.899 & 66980256 & 1649430/2.920 & 66980256 & 1351088/3.564 & 66980256 \\
093 & 15876/4.674 & 32768 & 15876/4.674 & 32768 & 15876/4.674 & 32768 \\
094 & 24640/3.851 & 318528 & 21573/4.399 & 424800 & 20380/4.656 & 424800 \\
095 & 420852/4.954 & 0 & 420852/4.954 & 0 & 420852/4.954 & 0 \\
096 & 26984/3.120 & 79488 & 26984/3.120 & 79488 & 26984/3.120 & 79488 \\
097 & 2102348/5.466 & 21311488 & 2102348/5.466 & 21311488 & 2102250/5.466 & 21311488 \\
098 & 89825/1.975 & 0 & 89825/1.975 & 0 & 89825/1.975 & 0 \\
099 & 339475/4.781 & 10698752 & 339475/4.781 & 10698752 & 329676/4.923 & 10698752 \\
100 & 57811/2.616 & 72192 & 57811/2.616 & 72192 & 57811/2.616 & 72192 \\
\end{longtable}

\endgroup

\section{Cache同计划完整配对}\label{app:cache}
表\ref{tab:cache-detail-1}至表\ref{tab:cache-detail-5}共列出500组配对。
$T_B=F_B(X_B^*)$，$T_C(X_B)=F_C(X_B^*)$，
$T_C(X_C)=F_C(X_C^*)$；三比值由式\eqref{eq:cache-factors}计算。
表中同时列出$C(X_B)$与$C(X_C)$各自的字节命中率；硬件效应图和同计划工期解释只使用$C(X_B)$字段。
每组的$C(X_B)$沿用该组$B$的最终计划，因而三段工期的计划身份保持一致。
\begingroup
\scriptsize
\setlength{\tabcolsep}{2pt}
\setlength{\LTleft}{\fill}\setlength{\LTright}{\fill}
\renewcommand{\arraystretch}{1.03}
\begin{longtable}{crrrrrrrr}
\caption{1核Cache配对（100图）}\label{tab:cache-detail-1}\\
\toprule
图 & $T_B$ & $T_C(X_B)$ & $T_C(X_C)$ & $S_{hw}$ & $S_{ad}$ & $S_{opt}$ & $C(X_B)$命中/\% & $C(X_C)$命中/\% \\
\midrule
\endfirsthead
\multicolumn{9}{c}{续表}\\
\toprule
图 & $T_B$ & $T_C(X_B)$ & $T_C(X_C)$ & $S_{hw}$ & $S_{ad}$ & $S_{opt}$ & $C(X_B)$命中/\% & $C(X_C)$命中/\% \\
\midrule
\endhead
\midrule\multicolumn{9}{r}{续下页}\\
\endfoot
\bottomrule
\endlastfoot
001 & 260797 & 257294 & 257294 & 1.0136 & 1.0000 & 1.0136 & 25.93 & 25.93 \\
002 & 269421 & 269421 & 267693 & 1.0000 & 1.0065 & 1.0065 & 27.74 & 0.00 \\
003 & 592981 & 592981 & 592981 & 1.0000 & 1.0000 & 1.0000 & 0.00 & 0.00 \\
004 & 415254 & 415254 & 402881 & 1.0000 & 1.0307 & 1.0307 & 0.00 & 10.41 \\
005 & 93910 & 93910 & 93910 & 1.0000 & 1.0000 & 1.0000 & 0.00 & 0.00 \\
006 & 76655 & 76655 & 76655 & 1.0000 & 1.0000 & 1.0000 & 1.45 & 1.45 \\
007 & 1210160 & 1210160 & 1155794 & 1.0000 & 1.0470 & 1.0470 & 0.00 & 21.72 \\
008 & 405761 & 405761 & 335225 & 1.0000 & 1.2104 & 1.2104 & 0.00 & 0.00 \\
009 & 160530 & 160530 & 160530 & 1.0000 & 1.0000 & 1.0000 & 0.00 & 0.00 \\
010 & 82266 & 82266 & 82266 & 1.0000 & 1.0000 & 1.0000 & 0.00 & 0.00 \\
011 & 186023 & 186023 & 186023 & 1.0000 & 1.0000 & 1.0000 & 0.65 & 0.65 \\
012 & 50002 & 49967 & 49967 & 1.0007 & 1.0000 & 1.0007 & 1.65 & 1.65 \\
013 & 379847 & 379847 & 309603 & 1.0000 & 1.2269 & 1.2269 & 0.00 & 0.00 \\
014 & 20060132 & 20036512 & 20036512 & 1.0012 & 1.0000 & 1.0012 & 5.46 & 5.46 \\
015 & 178946 & 178535 & 178342 & 1.0023 & 1.0011 & 1.0034 & 14.70 & 16.93 \\
016 & 7715523 & 7715523 & 7715523 & 1.0000 & 1.0000 & 1.0000 & 0.00 & 0.00 \\
017 & 167826 & 167609 & 167609 & 1.0013 & 1.0000 & 1.0013 & 10.37 & 10.37 \\
018 & 1050284 & 1048574 & 1003110 & 1.0016 & 1.0453 & 1.0470 & 19.95 & 34.38 \\
019 & 70264 & 70264 & 70264 & 1.0000 & 1.0000 & 1.0000 & 0.00 & 0.00 \\
020 & 585872 & 585872 & 542830 & 1.0000 & 1.0793 & 1.0793 & 0.00 & 0.00 \\
021 & 7425410 & 7383803 & 6899824 & 1.0056 & 1.0701 & 1.0762 & 8.35 & 17.19 \\
022 & 54906 & 54906 & 54906 & 1.0000 & 1.0000 & 1.0000 & 10.42 & 10.42 \\
023 & 116685 & 116571 & 116571 & 1.0010 & 1.0000 & 1.0010 & 9.88 & 9.88 \\
024 & 2555343 & 2555343 & 2555343 & 1.0000 & 1.0000 & 1.0000 & 0.00 & 0.00 \\
025 & 6233412 & 6233412 & 5823114 & 1.0000 & 1.0705 & 1.0705 & 0.00 & 0.00 \\
026 & 95307 & 95307 & 95307 & 1.0000 & 1.0000 & 1.0000 & 0.00 & 0.00 \\
027 & 2341632 & 2341632 & 2341632 & 1.0000 & 1.0000 & 1.0000 & 0.00 & 0.00 \\
028 & 43281233 & 42606375 & 42606375 & 1.0158 & 1.0000 & 1.0158 & 35.53 & 35.53 \\
029 & 37821 & 37581 & 37581 & 1.0064 & 1.0000 & 1.0064 & 14.03 & 14.03 \\
030 & 3462984 & 3461842 & 2940400 & 1.0003 & 1.1773 & 1.1777 & 6.11 & 14.59 \\
031 & 1508033 & 1505462 & 1505462 & 1.0017 & 1.0000 & 1.0017 & 7.95 & 7.95 \\
032 & 34533 & 34533 & 34533 & 1.0000 & 1.0000 & 1.0000 & 0.00 & 0.00 \\
033 & 2004121 & 2003734 & 1817490 & 1.0002 & 1.1025 & 1.1027 & 4.05 & 7.03 \\
034 & 927435 & 926166 & 842733 & 1.0014 & 1.0990 & 1.1005 & 11.63 & 30.90 \\
035 & 175454 & 175454 & 174846 & 1.0000 & 1.0035 & 1.0035 & 11.81 & 11.81 \\
036 & 1208827 & 1208827 & 1075353 & 1.0000 & 1.1241 & 1.1241 & 0.00 & 0.00 \\
037 & 244641 & 244535 & 244535 & 1.0004 & 1.0000 & 1.0004 & 0.66 & 0.66 \\
038 & 1589087 & 1583771 & 1583771 & 1.0034 & 1.0000 & 1.0034 & 15.67 & 15.67 \\
039 & 1752064 & 1729915 & 1729915 & 1.0128 & 1.0000 & 1.0128 & 10.06 & 10.06 \\
040 & 254945 & 254840 & 254840 & 1.0004 & 1.0000 & 1.0004 & 0.79 & 0.79 \\
041 & 7050742 & 7039953 & 6251871 & 1.0015 & 1.1261 & 1.1278 & 4.85 & 7.44 \\
042 & 1771204 & 1769685 & 1769685 & 1.0009 & 1.0000 & 1.0009 & 10.93 & 10.93 \\
043 & 675739 & 675739 & 675739 & 1.0000 & 1.0000 & 1.0000 & 3.29 & 3.29 \\
044 & 92892 & 92892 & 92892 & 1.0000 & 1.0000 & 1.0000 & 0.00 & 0.00 \\
045 & 135584 & 133243 & 126392 & 1.0176 & 1.0542 & 1.0727 & 13.57 & 9.15 \\
046 & 227889 & 227889 & 227889 & 1.0000 & 1.0000 & 1.0000 & 31.73 & 31.73 \\
047 & 324985 & 324985 & 324985 & 1.0000 & 1.0000 & 1.0000 & 0.00 & 0.00 \\
048 & 223338 & 223338 & 222614 & 1.0000 & 1.0033 & 1.0033 & 0.00 & 0.00 \\
049 & 199842 & 199842 & 199842 & 1.0000 & 1.0000 & 1.0000 & 0.00 & 0.00 \\
050 & 149846 & 149846 & 149470 & 1.0000 & 1.0025 & 1.0025 & 14.62 & 14.62 \\
051 & 607628 & 607628 & 607628 & 1.0000 & 1.0000 & 1.0000 & 0.00 & 0.00 \\
052 & 178673 & 178673 & 178673 & 1.0000 & 1.0000 & 1.0000 & 4.16 & 4.16 \\
053 & 1448004 & 1426552 & 1426552 & 1.0150 & 1.0000 & 1.0150 & 39.64 & 39.64 \\
054 & 796718 & 796718 & 796718 & 1.0000 & 1.0000 & 1.0000 & 0.00 & 0.00 \\
055 & 121890 & 120648 & 120648 & 1.0103 & 1.0000 & 1.0103 & 9.72 & 9.72 \\
056 & 278736 & 278736 & 277966 & 1.0000 & 1.0028 & 1.0028 & 0.00 & 0.00 \\
057 & 58953 & 58106 & 58106 & 1.0146 & 1.0000 & 1.0146 & 20.05 & 20.05 \\
058 & 6441268 & 6441268 & 6441268 & 1.0000 & 1.0000 & 1.0000 & 0.00 & 0.00 \\
059 & 768400 & 768400 & 768400 & 1.0000 & 1.0000 & 1.0000 & 0.78 & 0.78 \\
060 & 1095459 & 1093395 & 1086432 & 1.0019 & 1.0064 & 1.0083 & 6.92 & 10.97 \\
061 & 88220 & 88220 & 81014 & 1.0000 & 1.0889 & 1.0889 & 0.00 & 0.00 \\
062 & 3341379 & 3341379 & 3151609 & 1.0000 & 1.0602 & 1.0602 & 0.00 & 0.00 \\
063 & 1213468 & 1213468 & 1147397 & 1.0000 & 1.0576 & 1.0576 & 0.00 & 0.00 \\
064 & 19235 & 19235 & 19235 & 1.0000 & 1.0000 & 1.0000 & 0.00 & 0.00 \\
065 & 42821 & 42821 & 42821 & 1.0000 & 1.0000 & 1.0000 & 0.00 & 0.00 \\
066 & 358742 & 358742 & 358742 & 1.0000 & 1.0000 & 1.0000 & 0.00 & 0.00 \\
067 & 63782329 & 63620658 & 63620658 & 1.0025 & 1.0000 & 1.0025 & 7.98 & 7.98 \\
068 & 360987 & 360987 & 360987 & 1.0000 & 1.0000 & 1.0000 & 0.07 & 0.07 \\
069 & 23559 & 23559 & 23559 & 1.0000 & 1.0000 & 1.0000 & 0.00 & 0.00 \\
070 & 123146 & 123146 & 123146 & 1.0000 & 1.0000 & 1.0000 & 0.00 & 0.00 \\
071 & 18622 & 18622 & 18337 & 1.0000 & 1.0155 & 1.0155 & 0.00 & 0.00 \\
072 & 24907218 & 24884286 & 24884286 & 1.0009 & 1.0000 & 1.0009 & 3.67 & 3.67 \\
073 & 10067403 & 10067403 & 10067403 & 1.0000 & 1.0000 & 1.0000 & 0.00 & 0.00 \\
074 & 1737350 & 1731776 & 1731776 & 1.0032 & 1.0000 & 1.0032 & 13.55 & 13.55 \\
075 & 1433562 & 1433562 & 1433562 & 1.0000 & 1.0000 & 1.0000 & 0.00 & 0.00 \\
076 & 7839069 & 7831720 & 6603565 & 1.0009 & 1.1860 & 1.1871 & 4.95 & 8.72 \\
077 & 552271 & 552271 & 552271 & 1.0000 & 1.0000 & 1.0000 & 1.79 & 1.79 \\
078 & 2259940 & 2189192 & 2189192 & 1.0323 & 1.0000 & 1.0323 & 71.88 & 71.88 \\
079 & 35817831 & 35637258 & 35637258 & 1.0051 & 1.0000 & 1.0051 & 7.61 & 7.61 \\
080 & 329200 & 329096 & 329096 & 1.0003 & 1.0000 & 1.0003 & 0.21 & 0.21 \\
081 & 154087 & 151865 & 151865 & 1.0146 & 1.0000 & 1.0146 & 26.26 & 26.26 \\
082 & 717008 & 717008 & 717008 & 1.0000 & 1.0000 & 1.0000 & 0.00 & 0.00 \\
083 & 1040648 & 930103 & 930103 & 1.1189 & 1.0000 & 1.1189 & 76.91 & 76.91 \\
084 & 1762609 & 1762609 & 1552121 & 1.0000 & 1.1356 & 1.1356 & 0.00 & 0.00 \\
085 & 3187417 & 3187417 & 3187417 & 1.0000 & 1.0000 & 1.0000 & 0.00 & 0.00 \\
086 & 114843 & 114843 & 114843 & 1.0000 & 1.0000 & 1.0000 & 0.00 & 0.00 \\
087 & 3971407 & 3967115 & 3967115 & 1.0011 & 1.0000 & 1.0011 & 18.11 & 18.11 \\
088 & 240293 & 240293 & 240293 & 1.0000 & 1.0000 & 1.0000 & 0.00 & 0.00 \\
089 & 694728 & 694728 & 694728 & 1.0000 & 1.0000 & 1.0000 & 2.11 & 2.11 \\
090 & 1090066 & 998222 & 998222 & 1.0920 & 1.0000 & 1.0920 & 61.68 & 61.68 \\
091 & 10669316 & 10635127 & 9599880 & 1.0032 & 1.1078 & 1.1114 & 5.39 & 8.13 \\
092 & 4203197 & 4076456 & 4076456 & 1.0311 & 1.0000 & 1.0311 & 66.06 & 66.06 \\
093 & 85150 & 85150 & 77781 & 1.0000 & 1.0947 & 1.0947 & 1.02 & 0.00 \\
094 & 92511 & 92511 & 92511 & 1.0000 & 1.0000 & 1.0000 & 4.46 & 4.46 \\
095 & 1651133 & 1651133 & 1316513 & 1.0000 & 1.2542 & 1.2542 & 0.00 & 0.00 \\
096 & 80228 & 80228 & 80228 & 1.0000 & 1.0000 & 1.0000 & 0.00 & 0.00 \\
097 & 11385302 & 11381681 & 11365969 & 1.0003 & 1.0014 & 1.0017 & 4.63 & 11.97 \\
098 & 174782 & 174782 & 174782 & 1.0000 & 1.0000 & 1.0000 & 2.72 & 2.72 \\
099 & 1754664 & 1750575 & 1750575 & 1.0023 & 1.0000 & 1.0023 & 13.95 & 13.95 \\
100 & 145067 & 145067 & 145067 & 1.0000 & 1.0000 & 1.0000 & 0.00 & 0.00 \\
\end{longtable}

\begin{longtable}{crrrrrrrr}
\caption{2核Cache配对（100图）}\label{tab:cache-detail-2}\\
\toprule
图 & $T_B$ & $T_C(X_B)$ & $T_C(X_C)$ & $S_{hw}$ & $S_{ad}$ & $S_{opt}$ & $C(X_B)$命中/\% & $C(X_C)$命中/\% \\
\midrule
\endfirsthead
\multicolumn{9}{c}{续表}\\
\toprule
图 & $T_B$ & $T_C(X_B)$ & $T_C(X_C)$ & $S_{hw}$ & $S_{ad}$ & $S_{opt}$ & $C(X_B)$命中/\% & $C(X_C)$命中/\% \\
\midrule
\endhead
\midrule\multicolumn{9}{r}{续下页}\\
\endfoot
\bottomrule
\endlastfoot
001 & 116868 & 116868 & 116868 & 1.0000 & 1.0000 & 1.0000 & 0.00 & 0.00 \\
002 & 208461 & 205287 & 205287 & 1.0155 & 1.0000 & 1.0155 & 13.81 & 13.81 \\
003 & 562101 & 562101 & 562101 & 1.0000 & 1.0000 & 1.0000 & 5.21 & 5.21 \\
004 & 199074 & 199074 & 199074 & 1.0000 & 1.0000 & 1.0000 & 0.00 & 0.00 \\
005 & 98089 & 98089 & 98089 & 1.0000 & 1.0000 & 1.0000 & 0.46 & 0.46 \\
006 & 42696 & 42696 & 42696 & 1.0000 & 1.0000 & 1.0000 & 9.37 & 9.37 \\
007 & 507494 & 507494 & 507494 & 1.0000 & 1.0000 & 1.0000 & 0.00 & 0.00 \\
008 & 244266 & 244266 & 244266 & 1.0000 & 1.0000 & 1.0000 & 0.00 & 0.00 \\
009 & 142414 & 142414 & 142414 & 1.0000 & 1.0000 & 1.0000 & 3.70 & 3.70 \\
010 & 46568 & 46568 & 46568 & 1.0000 & 1.0000 & 1.0000 & 11.19 & 11.19 \\
011 & 92646 & 92646 & 92646 & 1.0000 & 1.0000 & 1.0000 & 0.00 & 0.00 \\
012 & 30058 & 30058 & 30058 & 1.0000 & 1.0000 & 1.0000 & 5.72 & 5.72 \\
013 & 132198 & 132198 & 132198 & 1.0000 & 1.0000 & 1.0000 & 0.00 & 0.00 \\
014 & 8912981 & 8904766 & 8904766 & 1.0009 & 1.0000 & 1.0009 & 24.60 & 24.60 \\
015 & 90509 & 90509 & 90509 & 1.0000 & 1.0000 & 1.0000 & 18.44 & 18.44 \\
016 & 7730025 & 7728365 & 7728365 & 1.0002 & 1.0000 & 1.0002 & 33.33 & 33.33 \\
017 & 81772 & 81772 & 81772 & 1.0000 & 1.0000 & 1.0000 & 1.30 & 1.30 \\
018 & 504468 & 502571 & 502571 & 1.0038 & 1.0000 & 1.0038 & 50.57 & 50.57 \\
019 & 41382 & 41382 & 41382 & 1.0000 & 1.0000 & 1.0000 & 13.16 & 13.16 \\
020 & 266580 & 266580 & 266580 & 1.0000 & 1.0000 & 1.0000 & 0.00 & 0.00 \\
021 & 3194570 & 3193634 & 3193634 & 1.0003 & 1.0000 & 1.0003 & 21.78 & 21.78 \\
022 & 28324 & 28324 & 28324 & 1.0000 & 1.0000 & 1.0000 & 3.14 & 3.14 \\
023 & 61501 & 61501 & 61501 & 1.0000 & 1.0000 & 1.0000 & 16.27 & 16.27 \\
024 & 2562266 & 2562266 & 2562266 & 1.0000 & 1.0000 & 1.0000 & 36.36 & 36.36 \\
025 & 2431840 & 2431840 & 2431840 & 1.0000 & 1.0000 & 1.0000 & 0.00 & 0.00 \\
026 & 55057 & 55057 & 55057 & 1.0000 & 1.0000 & 1.0000 & 4.29 & 4.29 \\
027 & 1120962 & 1120962 & 1120962 & 1.0000 & 1.0000 & 1.0000 & 8.90 & 8.90 \\
028 & 23638628 & 23615145 & 23615145 & 1.0010 & 1.0000 & 1.0010 & 0.47 & 0.47 \\
029 & 19740 & 19740 & 19740 & 1.0000 & 1.0000 & 1.0000 & 0.00 & 0.00 \\
030 & 1403594 & 1403449 & 1403449 & 1.0001 & 1.0000 & 1.0001 & 46.34 & 46.34 \\
031 & 712596 & 711944 & 711944 & 1.0009 & 1.0000 & 1.0009 & 43.05 & 43.05 \\
032 & 21040 & 21040 & 21040 & 1.0000 & 1.0000 & 1.0000 & 2.18 & 2.18 \\
033 & 889697 & 887951 & 887951 & 1.0020 & 1.0000 & 1.0020 & 35.16 & 35.16 \\
034 & 421427 & 421244 & 421244 & 1.0004 & 1.0000 & 1.0004 & 45.77 & 45.77 \\
035 & 151830 & 151830 & 151830 & 1.0000 & 1.0000 & 1.0000 & 12.87 & 12.87 \\
036 & 466068 & 466068 & 466068 & 1.0000 & 1.0000 & 1.0000 & 0.00 & 0.00 \\
037 & 105374 & 105374 & 105374 & 1.0000 & 1.0000 & 1.0000 & 0.00 & 0.00 \\
038 & 674500 & 674382 & 674382 & 1.0002 & 1.0000 & 1.0002 & 27.88 & 27.88 \\
039 & 708561 & 705615 & 705615 & 1.0042 & 1.0000 & 1.0042 & 37.02 & 37.02 \\
040 & 137746 & 137720 & 137720 & 1.0002 & 1.0000 & 1.0002 & 5.57 & 5.57 \\
041 & 2990398 & 2988290 & 2988290 & 1.0007 & 1.0000 & 1.0007 & 22.25 & 22.25 \\
042 & 849076 & 845668 & 845668 & 1.0040 & 1.0000 & 1.0040 & 39.40 & 39.40 \\
043 & 353851 & 353851 & 353851 & 1.0000 & 1.0000 & 1.0000 & 26.16 & 26.16 \\
044 & 73924 & 73924 & 73924 & 1.0000 & 1.0000 & 1.0000 & 0.00 & 0.00 \\
045 & 63059 & 63059 & 63059 & 1.0000 & 1.0000 & 1.0000 & 0.00 & 0.00 \\
046 & 166578 & 137552 & 137552 & 1.2110 & 1.0000 & 1.2110 & 57.36 & 57.36 \\
047 & 325068 & 325068 & 325068 & 1.0000 & 1.0000 & 1.0000 & 0.01 & 0.01 \\
048 & 225562 & 225562 & 225562 & 1.0000 & 1.0000 & 1.0000 & 0.02 & 0.02 \\
049 & 178903 & 178903 & 178903 & 1.0000 & 1.0000 & 1.0000 & 0.23 & 0.23 \\
050 & 149158 & 149158 & 149158 & 1.0000 & 1.0000 & 1.0000 & 16.47 & 16.47 \\
051 & 609681 & 609681 & 609681 & 1.0000 & 1.0000 & 1.0000 & 45.83 & 45.83 \\
052 & 90903 & 90903 & 90903 & 1.0000 & 1.0000 & 1.0000 & 11.27 & 11.27 \\
053 & 828335 & 796880 & 796880 & 1.0395 & 1.0000 & 1.0395 & 47.43 & 47.43 \\
054 & 811945 & 811945 & 811945 & 1.0000 & 1.0000 & 1.0000 & 0.95 & 0.95 \\
055 & 61686 & 61686 & 61686 & 1.0000 & 1.0000 & 1.0000 & 10.61 & 10.61 \\
056 & 281905 & 281905 & 281905 & 1.0000 & 1.0000 & 1.0000 & 0.01 & 0.01 \\
057 & 30183 & 30183 & 30183 & 1.0000 & 1.0000 & 1.0000 & 6.21 & 6.21 \\
058 & 2271316 & 2270206 & 2270206 & 1.0005 & 1.0000 & 1.0005 & 9.92 & 9.92 \\
059 & 369318 & 369318 & 369318 & 1.0000 & 1.0000 & 1.0000 & 0.00 & 0.00 \\
060 & 443955 & 443733 & 443733 & 1.0005 & 1.0000 & 1.0005 & 43.17 & 43.17 \\
061 & 42478 & 42478 & 42478 & 1.0000 & 1.0000 & 1.0000 & 5.97 & 5.97 \\
062 & 2610093 & 2525747 & 2525747 & 1.0334 & 1.0000 & 1.0334 & 8.45 & 8.45 \\
063 & 945557 & 914171 & 914171 & 1.0343 & 1.0000 & 1.0343 & 8.61 & 8.61 \\
064 & 20824 & 20824 & 20824 & 1.0000 & 1.0000 & 1.0000 & 0.43 & 0.43 \\
065 & 26610 & 26610 & 26610 & 1.0000 & 1.0000 & 1.0000 & 4.45 & 4.45 \\
066 & 342680 & 342680 & 342680 & 1.0000 & 1.0000 & 1.0000 & 3.56 & 3.56 \\
067 & 33768704 & 33768704 & 33768704 & 1.0000 & 1.0000 & 1.0000 & 0.00 & 0.00 \\
068 & 339779 & 339383 & 339318 & 1.0012 & 1.0002 & 1.0014 & 11.31 & 9.11 \\
069 & 25453 & 25453 & 25453 & 1.0000 & 1.0000 & 1.0000 & 3.92 & 3.92 \\
070 & 60121 & 60121 & 60121 & 1.0000 & 1.0000 & 1.0000 & 23.58 & 23.58 \\
071 & 19280 & 19280 & 19280 & 1.0000 & 1.0000 & 1.0000 & 3.50 & 3.50 \\
072 & 11872949 & 11862815 & 11862815 & 1.0009 & 1.0000 & 1.0009 & 12.17 & 12.17 \\
073 & 5980423 & 5980423 & 5980423 & 1.0000 & 1.0000 & 1.0000 & 0.00 & 0.00 \\
074 & 839999 & 837941 & 837941 & 1.0025 & 1.0000 & 1.0025 & 29.28 & 29.28 \\
075 & 1440267 & 1440267 & 1440267 & 1.0000 & 1.0000 & 1.0000 & 0.01 & 0.01 \\
076 & 3065214 & 3063843 & 3063843 & 1.0004 & 1.0000 & 1.0004 & 27.88 & 27.88 \\
077 & 248734 & 248734 & 248734 & 1.0000 & 1.0000 & 1.0000 & 2.61 & 2.61 \\
078 & 1260358 & 1260358 & 1260358 & 1.0000 & 1.0000 & 1.0000 & 0.00 & 0.00 \\
079 & 15446486 & 15446486 & 15446486 & 1.0000 & 1.0000 & 1.0000 & 0.00 & 0.00 \\
080 & 145926 & 145674 & 145674 & 1.0017 & 1.0000 & 1.0017 & 43.60 & 43.60 \\
081 & 76860 & 76860 & 76860 & 1.0000 & 1.0000 & 1.0000 & 0.00 & 0.00 \\
082 & 720881 & 720868 & 720868 & 1.0000 & 1.0000 & 1.0000 & 5.64 & 5.64 \\
083 & 775372 & 755398 & 755398 & 1.0264 & 1.0000 & 1.0264 & 43.21 & 43.21 \\
084 & 1255052 & 1255052 & 1255052 & 1.0000 & 1.0000 & 1.0000 & 0.00 & 0.00 \\
085 & 3207447 & 3207447 & 3207447 & 1.0000 & 1.0000 & 1.0000 & 7.18 & 7.18 \\
086 & 119261 & 119261 & 119261 & 1.0000 & 1.0000 & 1.0000 & 0.44 & 0.44 \\
087 & 2016479 & 2013301 & 2013301 & 1.0016 & 1.0000 & 1.0016 & 29.43 & 29.43 \\
088 & 243623 & 243623 & 243623 & 1.0000 & 1.0000 & 1.0000 & 0.00 & 0.00 \\
089 & 346832 & 346832 & 346832 & 1.0000 & 1.0000 & 1.0000 & 0.00 & 0.00 \\
090 & 836229 & 766154 & 766154 & 1.0915 & 1.0000 & 1.0915 & 31.51 & 31.51 \\
091 & 4571759 & 4569395 & 4569395 & 1.0005 & 1.0000 & 1.0005 & 17.99 & 17.99 \\
092 & 2828252 & 2358868 & 2358868 & 1.1990 & 1.0000 & 1.1990 & 75.92 & 75.92 \\
093 & 37827 & 37827 & 37827 & 1.0000 & 1.0000 & 1.0000 & 0.00 & 0.00 \\
094 & 48073 & 48073 & 48073 & 1.0000 & 1.0000 & 1.0000 & 6.43 & 6.43 \\
095 & 1042890 & 1042890 & 1042890 & 1.0000 & 1.0000 & 1.0000 & 0.00 & 0.00 \\
096 & 44010 & 44010 & 44010 & 1.0000 & 1.0000 & 1.0000 & 24.27 & 24.27 \\
097 & 5617602 & 5617602 & 5617602 & 1.0000 & 1.0000 & 1.0000 & 0.00 & 0.00 \\
098 & 90360 & 90360 & 90360 & 1.0000 & 1.0000 & 1.0000 & 0.00 & 0.00 \\
099 & 808916 & 806053 & 806053 & 1.0036 & 1.0000 & 1.0036 & 43.77 & 43.77 \\
100 & 81080 & 81080 & 81080 & 1.0000 & 1.0000 & 1.0000 & 3.73 & 3.73 \\
\end{longtable}

\begin{longtable}{crrrrrrrr}
\caption{3核Cache配对（100图）}\label{tab:cache-detail-3}\\
\toprule
图 & $T_B$ & $T_C(X_B)$ & $T_C(X_C)$ & $S_{hw}$ & $S_{ad}$ & $S_{opt}$ & $C(X_B)$命中/\% & $C(X_C)$命中/\% \\
\midrule
\endfirsthead
\multicolumn{9}{c}{续表}\\
\toprule
图 & $T_B$ & $T_C(X_B)$ & $T_C(X_C)$ & $S_{hw}$ & $S_{ad}$ & $S_{opt}$ & $C(X_B)$命中/\% & $C(X_C)$命中/\% \\
\midrule
\endhead
\midrule\multicolumn{9}{r}{续下页}\\
\endfoot
\bottomrule
\endlastfoot
001 & 78545 & 78545 & 78545 & 1.0000 & 1.0000 & 1.0000 & 0.00 & 0.00 \\
002 & 261945 & 261945 & 261945 & 1.0000 & 1.0000 & 1.0000 & 0.00 & 0.00 \\
003 & 580517 & 580515 & 580515 & 1.0000 & 1.0000 & 1.0000 & 1.54 & 1.54 \\
004 & 134043 & 134043 & 134043 & 1.0000 & 1.0000 & 1.0000 & 0.00 & 0.00 \\
005 & 102393 & 102392 & 102392 & 1.0000 & 1.0000 & 1.0000 & 3.00 & 3.00 \\
006 & 31152 & 31152 & 31152 & 1.0000 & 1.0000 & 1.0000 & 5.25 & 5.25 \\
007 & 338094 & 338094 & 338094 & 1.0000 & 1.0000 & 1.0000 & 0.00 & 0.00 \\
008 & 163359 & 163359 & 163359 & 1.0000 & 1.0000 & 1.0000 & 0.00 & 0.00 \\
009 & 153261 & 153261 & 153261 & 1.0000 & 1.0000 & 1.0000 & 0.00 & 0.00 \\
010 & 32889 & 32889 & 32889 & 1.0000 & 1.0000 & 1.0000 & 15.25 & 15.25 \\
011 & 69625 & 69625 & 69625 & 1.0000 & 1.0000 & 1.0000 & 56.69 & 56.69 \\
012 & 20661 & 20661 & 20661 & 1.0000 & 1.0000 & 1.0000 & 7.49 & 7.49 \\
013 & 88351 & 88351 & 88351 & 1.0000 & 1.0000 & 1.0000 & 0.00 & 0.00 \\
014 & 5993652 & 5957399 & 5957399 & 1.0061 & 1.0000 & 1.0061 & 28.96 & 28.96 \\
015 & 62280 & 62280 & 62280 & 1.0000 & 1.0000 & 1.0000 & 22.20 & 22.20 \\
016 & 7729868 & 7728076 & 7728076 & 1.0002 & 1.0000 & 1.0002 & 27.27 & 27.27 \\
017 & 55389 & 55389 & 55389 & 1.0000 & 1.0000 & 1.0000 & 12.98 & 12.98 \\
018 & 341047 & 338488 & 338488 & 1.0076 & 1.0000 & 1.0076 & 62.56 & 62.56 \\
019 & 28657 & 28657 & 28657 & 1.0000 & 1.0000 & 1.0000 & 23.26 & 23.26 \\
020 & 177950 & 177950 & 177950 & 1.0000 & 1.0000 & 1.0000 & 0.00 & 0.00 \\
021 & 2144024 & 2142851 & 2142851 & 1.0005 & 1.0000 & 1.0005 & 54.59 & 54.59 \\
022 & 19026 & 19026 & 19026 & 1.0000 & 1.0000 & 1.0000 & 6.08 & 6.08 \\
023 & 41403 & 41403 & 41403 & 1.0000 & 1.0000 & 1.0000 & 10.54 & 10.54 \\
024 & 2569688 & 2567896 & 2567896 & 1.0007 & 1.0000 & 1.0007 & 27.27 & 27.27 \\
025 & 1622164 & 1622164 & 1622164 & 1.0000 & 1.0000 & 1.0000 & 0.00 & 0.00 \\
026 & 41665 & 41665 & 41665 & 1.0000 & 1.0000 & 1.0000 & 16.37 & 16.37 \\
027 & 2077652 & 2077652 & 2077652 & 1.0000 & 1.0000 & 1.0000 & 0.24 & 0.24 \\
028 & 16770711 & 16770711 & 16770711 & 1.0000 & 1.0000 & 1.0000 & 0.00 & 0.00 \\
029 & 12743 & 12743 & 12743 & 1.0000 & 1.0000 & 1.0000 & 0.00 & 0.00 \\
030 & 933857 & 933236 & 933236 & 1.0007 & 1.0000 & 1.0007 & 54.65 & 54.65 \\
031 & 473822 & 473000 & 473000 & 1.0017 & 1.0000 & 1.0017 & 49.73 & 49.73 \\
032 & 14175 & 14175 & 14175 & 1.0000 & 1.0000 & 1.0000 & 4.26 & 4.26 \\
033 & 592065 & 590152 & 590152 & 1.0032 & 1.0000 & 1.0032 & 38.26 & 38.26 \\
034 & 276994 & 276560 & 276560 & 1.0016 & 1.0000 & 1.0016 & 54.75 & 54.75 \\
035 & 155604 & 155561 & 155561 & 1.0003 & 1.0000 & 1.0003 & 25.87 & 25.87 \\
036 & 311345 & 311345 & 311345 & 1.0000 & 1.0000 & 1.0000 & 0.00 & 0.00 \\
037 & 74038 & 74038 & 74038 & 1.0000 & 1.0000 & 1.0000 & 58.97 & 58.97 \\
038 & 452248 & 452022 & 452022 & 1.0005 & 1.0000 & 1.0005 & 35.73 & 35.73 \\
039 & 485565 & 483226 & 483226 & 1.0048 & 1.0000 & 1.0048 & 38.60 & 38.60 \\
040 & 110105 & 110105 & 110105 & 1.0000 & 1.0000 & 1.0000 & 0.00 & 0.00 \\
041 & 1991911 & 1989953 & 1989953 & 1.0010 & 1.0000 & 1.0010 & 26.05 & 26.05 \\
042 & 562868 & 561628 & 561628 & 1.0022 & 1.0000 & 1.0022 & 43.33 & 43.33 \\
043 & 268777 & 268777 & 268777 & 1.0000 & 1.0000 & 1.0000 & 8.99 & 8.99 \\
044 & 96051 & 96051 & 96051 & 1.0000 & 1.0000 & 1.0000 & 1.90 & 1.90 \\
045 & 42601 & 42601 & 42601 & 1.0000 & 1.0000 & 1.0000 & 0.00 & 0.00 \\
046 & 142176 & 121494 & 121494 & 1.1702 & 1.0000 & 1.1702 & 43.06 & 43.06 \\
047 & 328135 & 328133 & 328133 & 1.0000 & 1.0000 & 1.0000 & 4.51 & 4.51 \\
048 & 227793 & 227793 & 227793 & 1.0000 & 1.0000 & 1.0000 & 0.03 & 0.03 \\
049 & 187317 & 187317 & 187317 & 1.0000 & 1.0000 & 1.0000 & 1.33 & 1.33 \\
050 & 153757 & 153757 & 153757 & 1.0000 & 1.0000 & 1.0000 & 6.06 & 6.06 \\
051 & 611232 & 611232 & 611232 & 1.0000 & 1.0000 & 1.0000 & 61.11 & 61.11 \\
052 & 61071 & 61071 & 61071 & 1.0000 & 1.0000 & 1.0000 & 18.28 & 18.28 \\
053 & 1472686 & 1428109 & 1428109 & 1.0312 & 1.0000 & 1.0312 & 62.49 & 62.49 \\
054 & 806897 & 806897 & 806897 & 1.0000 & 1.0000 & 1.0000 & 0.93 & 0.93 \\
055 & 42047 & 42047 & 42047 & 1.0000 & 1.0000 & 1.0000 & 15.61 & 15.61 \\
056 & 285038 & 285038 & 285038 & 1.0000 & 1.0000 & 1.0000 & 0.16 & 0.16 \\
057 & 20705 & 20705 & 20705 & 1.0000 & 1.0000 & 1.0000 & 0.08 & 0.08 \\
058 & 1554831 & 1554831 & 1554831 & 1.0000 & 1.0000 & 1.0000 & 0.00 & 0.00 \\
059 & 266217 & 262333 & 262333 & 1.0148 & 1.0000 & 1.0148 & 68.32 & 68.32 \\
060 & 298210 & 297398 & 297398 & 1.0027 & 1.0000 & 1.0027 & 48.04 & 48.04 \\
061 & 30303 & 30303 & 30303 & 1.0000 & 1.0000 & 1.0000 & 6.22 & 6.22 \\
062 & 2854809 & 2854809 & 2854809 & 1.0000 & 1.0000 & 1.0000 & 0.00 & 0.00 \\
063 & 1070397 & 1070397 & 1070397 & 1.0000 & 1.0000 & 1.0000 & 0.00 & 0.00 \\
064 & 22179 & 22177 & 22177 & 1.0001 & 1.0000 & 1.0001 & 3.92 & 3.92 \\
065 & 18032 & 18032 & 18032 & 1.0000 & 1.0000 & 1.0000 & 4.31 & 4.31 \\
066 & 316559 & 316556 & 316556 & 1.0000 & 1.0000 & 1.0000 & 5.16 & 5.16 \\
067 & 23568859 & 23568859 & 23568859 & 1.0000 & 1.0000 & 1.0000 & 0.00 & 0.00 \\
068 & 362095 & 362095 & 362095 & 1.0000 & 1.0000 & 1.0000 & 0.00 & 0.00 \\
069 & 29048 & 29048 & 26284 & 1.0000 & 1.1052 & 1.1052 & 7.22 & 6.25 \\
070 & 39333 & 39333 & 39333 & 1.0000 & 1.0000 & 1.0000 & 27.30 & 27.30 \\
071 & 23372 & 23363 & 23363 & 1.0004 & 1.0000 & 1.0004 & 13.60 & 13.60 \\
072 & 7935290 & 7922729 & 7922729 & 1.0016 & 1.0000 & 1.0016 & 15.41 & 15.41 \\
073 & 4515227 & 4515227 & 4515227 & 1.0000 & 1.0000 & 1.0000 & 0.00 & 0.00 \\
074 & 563098 & 557771 & 557771 & 1.0096 & 1.0000 & 1.0096 & 33.41 & 33.41 \\
075 & 1442868 & 1442868 & 1442868 & 1.0000 & 1.0000 & 1.0000 & 0.01 & 0.01 \\
076 & 2050028 & 2047868 & 2047868 & 1.0011 & 1.0000 & 1.0011 & 33.69 & 33.69 \\
077 & 168976 & 168976 & 168976 & 1.0000 & 1.0000 & 1.0000 & 2.54 & 2.54 \\
078 & 1377249 & 1366507 & 1366507 & 1.0079 & 1.0000 & 1.0079 & 23.27 & 23.27 \\
079 & 10181692 & 10175315 & 10175315 & 1.0006 & 1.0000 & 1.0006 & 13.73 & 13.73 \\
080 & 100325 & 100325 & 100325 & 1.0000 & 1.0000 & 1.0000 & 0.00 & 0.00 \\
081 & 51943 & 51943 & 51943 & 1.0000 & 1.0000 & 1.0000 & 0.00 & 0.00 \\
082 & 691651 & 691651 & 691651 & 1.0000 & 1.0000 & 1.0000 & 5.69 & 5.69 \\
083 & 512623 & 402691 & 402691 & 1.2730 & 1.0000 & 1.2730 & 65.90 & 65.90 \\
084 & 836220 & 836220 & 836220 & 1.0000 & 1.0000 & 1.0000 & 0.00 & 0.00 \\
085 & 3130505 & 3130505 & 3130505 & 1.0000 & 1.0000 & 1.0000 & 0.00 & 0.00 \\
086 & 118403 & 118403 & 118403 & 1.0000 & 1.0000 & 1.0000 & 0.00 & 0.00 \\
087 & 1374901 & 1372451 & 1372451 & 1.0018 & 1.0000 & 1.0018 & 26.15 & 26.15 \\
088 & 243623 & 243623 & 243623 & 1.0000 & 1.0000 & 1.0000 & 0.00 & 0.00 \\
089 & 243614 & 243614 & 243614 & 1.0000 & 1.0000 & 1.0000 & 0.00 & 0.00 \\
090 & 761694 & 706360 & 706360 & 1.0783 & 1.0000 & 1.0783 & 28.06 & 28.06 \\
091 & 3090094 & 3070819 & 3070819 & 1.0063 & 1.0000 & 1.0063 & 24.46 & 24.46 \\
092 & 2231804 & 1792534 & 1792534 & 1.2451 & 1.0000 & 1.2451 & 65.42 & 65.42 \\
093 & 25103 & 25103 & 25103 & 1.0000 & 1.0000 & 1.0000 & 0.00 & 0.00 \\
094 & 32222 & 32222 & 32222 & 1.0000 & 1.0000 & 1.0000 & 31.08 & 31.08 \\
095 & 695775 & 695775 & 695775 & 1.0000 & 1.0000 & 1.0000 & 0.00 & 0.00 \\
096 & 30318 & 30318 & 30318 & 1.0000 & 1.0000 & 1.0000 & 27.43 & 27.43 \\
097 & 10246000 & 10246000 & 10246000 & 1.0000 & 1.0000 & 1.0000 & 0.04 & 0.04 \\
098 & 89825 & 89825 & 89825 & 1.0000 & 1.0000 & 1.0000 & 0.00 & 0.00 \\
099 & 544403 & 540586 & 540586 & 1.0071 & 1.0000 & 1.0071 & 51.04 & 51.04 \\
100 & 57811 & 57811 & 57811 & 1.0000 & 1.0000 & 1.0000 & 10.58 & 10.58 \\
\end{longtable}

\begin{longtable}{crrrrrrrr}
\caption{4核Cache配对（100图）}\label{tab:cache-detail-4}\\
\toprule
图 & $T_B$ & $T_C(X_B)$ & $T_C(X_C)$ & $S_{hw}$ & $S_{ad}$ & $S_{opt}$ & $C(X_B)$命中/\% & $C(X_C)$命中/\% \\
\midrule
\endfirsthead
\multicolumn{9}{c}{续表}\\
\toprule
图 & $T_B$ & $T_C(X_B)$ & $T_C(X_C)$ & $S_{hw}$ & $S_{ad}$ & $S_{opt}$ & $C(X_B)$命中/\% & $C(X_C)$命中/\% \\
\midrule
\endhead
\midrule\multicolumn{9}{r}{续下页}\\
\endfoot
\bottomrule
\endlastfoot
001 & 58984 & 58984 & 58984 & 1.0000 & 1.0000 & 1.0000 & 0.00 & 0.00 \\
002 & 129903 & 124800 & 124800 & 1.0409 & 1.0000 & 1.0409 & 14.86 & 14.86 \\
003 & 602212 & 602212 & 602212 & 1.0000 & 1.0000 & 1.0000 & 0.00 & 0.00 \\
004 & 104188 & 104188 & 104188 & 1.0000 & 1.0000 & 1.0000 & 0.00 & 0.00 \\
005 & 98413 & 98300 & 98300 & 1.0011 & 1.0000 & 1.0011 & 5.33 & 5.33 \\
006 & 31428 & 31428 & 31428 & 1.0000 & 1.0000 & 1.0000 & 9.17 & 9.17 \\
007 & 254398 & 254398 & 254398 & 1.0000 & 1.0000 & 1.0000 & 0.00 & 0.00 \\
008 & 123060 & 123060 & 123060 & 1.0000 & 1.0000 & 1.0000 & 0.00 & 0.00 \\
009 & 153261 & 153261 & 153261 & 1.0000 & 1.0000 & 1.0000 & 0.00 & 0.00 \\
010 & 29918 & 29858 & 29858 & 1.0020 & 1.0000 & 1.0020 & 10.30 & 10.30 \\
011 & 46880 & 46880 & 46880 & 1.0000 & 1.0000 & 1.0000 & 0.00 & 0.00 \\
012 & 15863 & 15863 & 15863 & 1.0000 & 1.0000 & 1.0000 & 9.28 & 9.28 \\
013 & 66982 & 66982 & 66982 & 1.0000 & 1.0000 & 1.0000 & 0.00 & 0.00 \\
014 & 4551951 & 4526148 & 4526148 & 1.0057 & 1.0000 & 1.0057 & 32.37 & 32.37 \\
015 & 48727 & 48669 & 48669 & 1.0012 & 1.0000 & 1.0012 & 20.89 & 20.89 \\
016 & 7715523 & 7715523 & 7715523 & 1.0000 & 1.0000 & 1.0000 & 0.00 & 0.00 \\
017 & 41209 & 41209 & 41209 & 1.0000 & 1.0000 & 1.0000 & 31.85 & 31.85 \\
018 & 253676 & 253384 & 253384 & 1.0012 & 1.0000 & 1.0012 & 56.24 & 56.24 \\
019 & 24099 & 24099 & 24099 & 1.0000 & 1.0000 & 1.0000 & 19.75 & 19.75 \\
020 & 133704 & 133704 & 133704 & 1.0000 & 1.0000 & 1.0000 & 0.00 & 0.00 \\
021 & 1550324 & 1550324 & 1550324 & 1.0000 & 1.0000 & 1.0000 & 0.00 & 0.00 \\
022 & 14526 & 14526 & 14526 & 1.0000 & 1.0000 & 1.0000 & 6.08 & 6.08 \\
023 & 30142 & 30142 & 30142 & 1.0000 & 1.0000 & 1.0000 & 39.32 & 39.32 \\
024 & 2586270 & 2580375 & 2580375 & 1.0023 & 1.0000 & 1.0023 & 38.89 & 38.89 \\
025 & 1217444 & 1217444 & 1217444 & 1.0000 & 1.0000 & 1.0000 & 0.00 & 0.00 \\
026 & 42208 & 42208 & 42208 & 1.0000 & 1.0000 & 1.0000 & 28.13 & 28.13 \\
027 & 559196 & 559196 & 559196 & 1.0000 & 1.0000 & 1.0000 & 0.00 & 0.00 \\
028 & 13608312 & 13585418 & 13585418 & 1.0017 & 1.0000 & 1.0017 & 0.27 & 0.27 \\
029 & 11028 & 11028 & 11028 & 1.0000 & 1.0000 & 1.0000 & 0.00 & 0.00 \\
030 & 700342 & 700156 & 700156 & 1.0003 & 1.0000 & 1.0003 & 59.66 & 59.66 \\
031 & 357987 & 355839 & 355839 & 1.0060 & 1.0000 & 1.0060 & 48.80 & 48.80 \\
032 & 11004 & 11004 & 11004 & 1.0000 & 1.0000 & 1.0000 & 6.24 & 6.24 \\
033 & 443447 & 442667 & 442667 & 1.0018 & 1.0000 & 1.0018 & 42.67 & 42.67 \\
034 & 212170 & 212170 & 212170 & 1.0000 & 1.0000 & 1.0000 & 52.33 & 52.33 \\
035 & 169571 & 169571 & 169571 & 1.0000 & 1.0000 & 1.0000 & 0.00 & 0.00 \\
036 & 233584 & 233584 & 233584 & 1.0000 & 1.0000 & 1.0000 & 0.00 & 0.00 \\
037 & 53320 & 53320 & 53320 & 1.0000 & 1.0000 & 1.0000 & 0.00 & 0.00 \\
038 & 339624 & 337813 & 337813 & 1.0054 & 1.0000 & 1.0054 & 35.21 & 35.21 \\
039 & 395323 & 368055 & 368055 & 1.0741 & 1.0000 & 1.0741 & 47.82 & 47.82 \\
040 & 110202 & 110202 & 110105 & 1.0000 & 1.0009 & 1.0009 & 4.14 & 0.00 \\
041 & 1499144 & 1494347 & 1494347 & 1.0032 & 1.0000 & 1.0032 & 34.18 & 34.18 \\
042 & 426511 & 423652 & 423652 & 1.0067 & 1.0000 & 1.0067 & 43.14 & 43.14 \\
043 & 666701 & 666696 & 666696 & 1.0000 & 1.0000 & 1.0000 & 2.67 & 2.67 \\
044 & 90531 & 90531 & 90531 & 1.0000 & 1.0000 & 1.0000 & 2.70 & 2.70 \\
045 & 31878 & 31878 & 31878 & 1.0000 & 1.0000 & 1.0000 & 0.00 & 0.00 \\
046 & 179861 & 154683 & 149978 & 1.1628 & 1.0314 & 1.1992 & 60.40 & 43.02 \\
047 & 325683 & 325678 & 325678 & 1.0000 & 1.0000 & 1.0000 & 0.01 & 0.01 \\
048 & 227870 & 227866 & 227866 & 1.0000 & 1.0000 & 1.0000 & 0.00 & 0.00 \\
049 & 199561 & 199561 & 196356 & 1.0000 & 1.0163 & 1.0163 & 4.96 & 2.19 \\
050 & 149690 & 149690 & 149690 & 1.0000 & 1.0000 & 1.0000 & 0.00 & 0.00 \\
051 & 612783 & 612783 & 612783 & 1.0000 & 1.0000 & 1.0000 & 68.75 & 68.75 \\
052 & 45610 & 45610 & 45610 & 1.0000 & 1.0000 & 1.0000 & 16.59 & 16.59 \\
053 & 1472686 & 1428109 & 1428109 & 1.0312 & 1.0000 & 1.0312 & 62.49 & 62.49 \\
054 & 831731 & 831696 & 830580 & 1.0000 & 1.0013 & 1.0014 & 3.22 & 3.29 \\
055 & 32099 & 32099 & 32099 & 1.0000 & 1.0000 & 1.0000 & 21.72 & 21.72 \\
056 & 286916 & 286914 & 286914 & 1.0000 & 1.0000 & 1.0000 & 0.23 & 0.23 \\
057 & 16485 & 16485 & 16485 & 1.0000 & 1.0000 & 1.0000 & 7.74 & 7.74 \\
058 & 1154519 & 1142629 & 1142629 & 1.0104 & 1.0000 & 1.0104 & 26.90 & 26.90 \\
059 & 185216 & 185216 & 185216 & 1.0000 & 1.0000 & 1.0000 & 0.00 & 0.00 \\
060 & 221720 & 221694 & 221694 & 1.0001 & 1.0000 & 1.0001 & 48.17 & 48.17 \\
061 & 24024 & 24024 & 24024 & 1.0000 & 1.0000 & 1.0000 & 11.49 & 11.49 \\
062 & 2240290 & 2195265 & 2195265 & 1.0205 & 1.0000 & 1.0205 & 0.08 & 0.08 \\
063 & 742037 & 742037 & 742037 & 1.0000 & 1.0000 & 1.0000 & 1.74 & 1.74 \\
064 & 22068 & 22065 & 22065 & 1.0001 & 1.0000 & 1.0001 & 3.22 & 3.22 \\
065 & 13541 & 13541 & 13541 & 1.0000 & 1.0000 & 1.0000 & 10.92 & 10.92 \\
066 & 346934 & 346934 & 346934 & 1.0000 & 1.0000 & 1.0000 & 3.52 & 3.52 \\
067 & 18473155 & 18473155 & 18473155 & 1.0000 & 1.0000 & 1.0000 & 0.00 & 0.00 \\
068 & 351308 & 351305 & 351305 & 1.0000 & 1.0000 & 1.0000 & 3.29 & 3.29 \\
069 & 27749 & 27747 & 27745 & 1.0001 & 1.0001 & 1.0001 & 4.76 & 4.83 \\
070 & 31225 & 31192 & 31192 & 1.0011 & 1.0000 & 1.0011 & 39.93 & 39.93 \\
071 & 20350 & 20323 & 19251 & 1.0013 & 1.0557 & 1.0571 & 5.92 & 2.44 \\
072 & 6003776 & 5991003 & 5991003 & 1.0021 & 1.0000 & 1.0021 & 17.31 & 17.31 \\
073 & 3903787 & 3903787 & 3903787 & 1.0000 & 1.0000 & 1.0000 & 0.00 & 0.00 \\
074 & 418987 & 415633 & 415633 & 1.0081 & 1.0000 & 1.0081 & 32.23 & 32.23 \\
075 & 1430810 & 1430808 & 1430808 & 1.0000 & 1.0000 & 1.0000 & 0.01 & 0.01 \\
076 & 1540201 & 1536306 & 1536306 & 1.0025 & 1.0000 & 1.0025 & 41.71 & 41.71 \\
077 & 127832 & 127591 & 127591 & 1.0019 & 1.0000 & 1.0019 & 5.08 & 5.08 \\
078 & 690924 & 690924 & 690924 & 1.0000 & 1.0000 & 1.0000 & 0.00 & 0.00 \\
079 & 7635624 & 7635624 & 7635624 & 1.0000 & 1.0000 & 1.0000 & 0.00 & 0.00 \\
080 & 111314 & 90584 & 90584 & 1.2288 & 1.0000 & 1.2288 & 47.85 & 47.85 \\
081 & 39588 & 39588 & 39588 & 1.0000 & 1.0000 & 1.0000 & 0.00 & 0.00 \\
082 & 766658 & 766658 & 766658 & 1.0000 & 1.0000 & 1.0000 & 8.95 & 8.95 \\
083 & 450029 & 333151 & 333151 & 1.3508 & 1.0000 & 1.3508 & 64.26 & 64.26 \\
084 & 629028 & 629028 & 629028 & 1.0000 & 1.0000 & 1.0000 & 0.00 & 0.00 \\
085 & 3203847 & 3203760 & 3203760 & 1.0000 & 1.0000 & 1.0000 & 17.92 & 17.92 \\
086 & 127542 & 127542 & 126423 & 1.0000 & 1.0089 & 1.0089 & 3.97 & 3.81 \\
087 & 1046846 & 1043492 & 1043492 & 1.0032 & 1.0000 & 1.0032 & 29.95 & 29.95 \\
088 & 243623 & 243623 & 243623 & 1.0000 & 1.0000 & 1.0000 & 0.00 & 0.00 \\
089 & 178160 & 178160 & 178160 & 1.0000 & 1.0000 & 1.0000 & 0.00 & 0.00 \\
090 & 731084 & 674035 & 674035 & 1.0846 & 1.0000 & 1.0846 & 36.43 & 36.43 \\
091 & 2315897 & 2307637 & 2307637 & 1.0036 & 1.0000 & 1.0036 & 28.04 & 28.04 \\
092 & 1927728 & 1470508 & 1470508 & 1.3109 & 1.0000 & 1.3109 & 64.79 & 64.79 \\
093 & 19779 & 19779 & 19779 & 1.0000 & 1.0000 & 1.0000 & 0.00 & 0.00 \\
094 & 24640 & 24640 & 24640 & 1.0000 & 1.0000 & 1.0000 & 38.58 & 38.58 \\
095 & 524422 & 524422 & 524422 & 1.0000 & 1.0000 & 1.0000 & 0.00 & 0.00 \\
096 & 26984 & 26984 & 26984 & 1.0000 & 1.0000 & 1.0000 & 8.35 & 8.35 \\
097 & 2626908 & 2626908 & 2626908 & 1.0000 & 1.0000 & 1.0000 & 4.17 & 4.17 \\
098 & 89963 & 89963 & 89963 & 1.0000 & 1.0000 & 1.0000 & 0.00 & 0.00 \\
099 & 405336 & 403555 & 403555 & 1.0044 & 1.0000 & 1.0044 & 55.68 & 55.68 \\
100 & 57948 & 57948 & 57948 & 1.0000 & 1.0000 & 1.0000 & 24.88 & 24.88 \\
\end{longtable}

\begin{longtable}{crrrrrrrr}
\caption{5核Cache配对（100图）}\label{tab:cache-detail-5}\\
\toprule
图 & $T_B$ & $T_C(X_B)$ & $T_C(X_C)$ & $S_{hw}$ & $S_{ad}$ & $S_{opt}$ & $C(X_B)$命中/\% & $C(X_C)$命中/\% \\
\midrule
\endfirsthead
\multicolumn{9}{c}{续表}\\
\toprule
图 & $T_B$ & $T_C(X_B)$ & $T_C(X_C)$ & $S_{hw}$ & $S_{ad}$ & $S_{opt}$ & $C(X_B)$命中/\% & $C(X_C)$命中/\% \\
\midrule
\endhead
\midrule\multicolumn{9}{r}{续下页}\\
\endfoot
\bottomrule
\endlastfoot
001 & 47502 & 47502 & 47502 & 1.0000 & 1.0000 & 1.0000 & 0.00 & 0.00 \\
002 & 146574 & 139049 & 139049 & 1.0541 & 1.0000 & 1.0541 & 15.52 & 15.52 \\
003 & 541371 & 541371 & 541371 & 1.0000 & 1.0000 & 1.0000 & 1.59 & 1.59 \\
004 & 88341 & 88341 & 88341 & 1.0000 & 1.0000 & 1.0000 & 0.00 & 0.00 \\
005 & 95618 & 95618 & 95618 & 1.0000 & 1.0000 & 1.0000 & 0.00 & 0.00 \\
006 & 31753 & 31753 & 31753 & 1.0000 & 1.0000 & 1.0000 & 9.12 & 9.12 \\
007 & 203834 & 203834 & 203834 & 1.0000 & 1.0000 & 1.0000 & 0.00 & 0.00 \\
008 & 100603 & 100603 & 100603 & 1.0000 & 1.0000 & 1.0000 & 0.00 & 0.00 \\
009 & 139951 & 138757 & 138757 & 1.0086 & 1.0000 & 1.0086 & 15.66 & 15.66 \\
010 & 31019 & 30621 & 30621 & 1.0130 & 1.0000 & 1.0130 & 18.72 & 18.72 \\
011 & 47287 & 46777 & 46777 & 1.0109 & 1.0000 & 1.0109 & 62.31 & 62.31 \\
012 & 13406 & 13406 & 13406 & 1.0000 & 1.0000 & 1.0000 & 7.68 & 7.68 \\
013 & 54363 & 54363 & 54363 & 1.0000 & 1.0000 & 1.0000 & 0.00 & 0.00 \\
014 & 3712893 & 3700752 & 3700752 & 1.0033 & 1.0000 & 1.0033 & 29.39 & 29.39 \\
015 & 43665 & 42872 & 42872 & 1.0185 & 1.0000 & 1.0185 & 24.87 & 24.87 \\
016 & 7741232 & 7740666 & 7740666 & 1.0001 & 1.0000 & 1.0001 & 33.33 & 33.33 \\
017 & 34231 & 34231 & 34231 & 1.0000 & 1.0000 & 1.0000 & 21.62 & 21.62 \\
018 & 205219 & 201659 & 201659 & 1.0177 & 1.0000 & 1.0177 & 57.22 & 57.22 \\
019 & 24202 & 24202 & 24202 & 1.0000 & 1.0000 & 1.0000 & 1.99 & 1.99 \\
020 & 107650 & 107650 & 107650 & 1.0000 & 1.0000 & 1.0000 & 0.00 & 0.00 \\
021 & 1250168 & 1249380 & 1249380 & 1.0006 & 1.0000 & 1.0006 & 57.00 & 57.00 \\
022 & 11824 & 11824 & 11824 & 1.0000 & 1.0000 & 1.0000 & 8.85 & 8.85 \\
023 & 24517 & 24517 & 24517 & 1.0000 & 1.0000 & 1.0000 & 21.00 & 21.00 \\
024 & 2590252 & 2579548 & 2579548 & 1.0041 & 1.0000 & 1.0041 & 39.13 & 39.13 \\
025 & 975282 & 975282 & 975282 & 1.0000 & 1.0000 & 1.0000 & 0.00 & 0.00 \\
026 & 45246 & 44374 & 41655 & 1.0197 & 1.0653 & 1.0862 & 12.51 & 4.29 \\
027 & 558753 & 558753 & 558753 & 1.0000 & 1.0000 & 1.0000 & 4.44 & 4.44 \\
028 & 11491035 & 11450794 & 11450794 & 1.0035 & 1.0000 & 1.0035 & 1.06 & 1.06 \\
029 & 9383 & 9383 & 9383 & 1.0000 & 1.0000 & 1.0000 & 0.00 & 0.00 \\
030 & 561102 & 560897 & 560897 & 1.0004 & 1.0000 & 1.0004 & 59.51 & 59.51 \\
031 & 285693 & 284489 & 284489 & 1.0042 & 1.0000 & 1.0042 & 53.15 & 53.15 \\
032 & 10496 & 10479 & 10479 & 1.0016 & 1.0000 & 1.0016 & 4.26 & 4.26 \\
033 & 355102 & 354514 & 354514 & 1.0017 & 1.0000 & 1.0017 & 31.27 & 31.27 \\
034 & 170371 & 168925 & 168925 & 1.0086 & 1.0000 & 1.0086 & 58.25 & 58.25 \\
035 & 157857 & 157840 & 157337 & 1.0001 & 1.0032 & 1.0033 & 21.78 & 21.17 \\
036 & 187182 & 187182 & 187182 & 1.0000 & 1.0000 & 1.0000 & 0.00 & 0.00 \\
037 & 43800 & 43013 & 43013 & 1.0183 & 1.0000 & 1.0183 & 64.25 & 64.25 \\
038 & 275133 & 272533 & 272533 & 1.0095 & 1.0000 & 1.0095 & 42.22 & 42.22 \\
039 & 431161 & 340476 & 340476 & 1.2663 & 1.0000 & 1.2663 & 37.09 & 37.09 \\
040 & 110337 & 110337 & 110105 & 1.0000 & 1.0021 & 1.0021 & 4.07 & 0.00 \\
041 & 1201884 & 1194643 & 1194643 & 1.0061 & 1.0000 & 1.0061 & 32.66 & 32.66 \\
042 & 342520 & 339457 & 339457 & 1.0090 & 1.0000 & 1.0090 & 49.37 & 49.37 \\
043 & 677656 & 677656 & 677656 & 1.0000 & 1.0000 & 1.0000 & 3.32 & 3.32 \\
044 & 76858 & 75919 & 75919 & 1.0124 & 1.0000 & 1.0124 & 7.03 & 7.03 \\
045 & 26164 & 26164 & 26164 & 1.0000 & 1.0000 & 1.0000 & 0.00 & 0.00 \\
046 & 181490 & 180844 & 180844 & 1.0036 & 1.0000 & 1.0036 & 5.76 & 5.76 \\
047 & 316261 & 316261 & 316261 & 1.0000 & 1.0000 & 1.0000 & 0.05 & 0.05 \\
048 & 228173 & 228173 & 226918 & 1.0000 & 1.0055 & 1.0055 & 0.49 & 0.14 \\
049 & 195369 & 195368 & 195368 & 1.0000 & 1.0000 & 1.0000 & 2.16 & 2.16 \\
050 & 147954 & 147851 & 147851 & 1.0007 & 1.0000 & 1.0007 & 8.11 & 8.11 \\
051 & 627634 & 627502 & 627502 & 1.0002 & 1.0000 & 1.0002 & 42.39 & 42.39 \\
052 & 37387 & 37387 & 37387 & 1.0000 & 1.0000 & 1.0000 & 26.03 & 26.03 \\
053 & 1472686 & 1428109 & 1428109 & 1.0312 & 1.0000 & 1.0312 & 62.49 & 62.49 \\
054 & 822478 & 822451 & 822451 & 1.0000 & 1.0000 & 1.0000 & 4.48 & 4.48 \\
055 & 25635 & 25635 & 25635 & 1.0000 & 1.0000 & 1.0000 & 22.16 & 22.16 \\
056 & 290020 & 290020 & 286484 & 1.0000 & 1.0123 & 1.0123 & 1.02 & 1.20 \\
057 & 14028 & 14028 & 14028 & 1.0000 & 1.0000 & 1.0000 & 3.94 & 3.94 \\
058 & 1004819 & 1004819 & 1004819 & 1.0000 & 1.0000 & 1.0000 & 0.00 & 0.00 \\
059 & 150782 & 150736 & 150736 & 1.0003 & 1.0000 & 1.0003 & 73.79 & 73.79 \\
060 & 175702 & 175626 & 175626 & 1.0004 & 1.0000 & 1.0004 & 52.58 & 52.58 \\
061 & 22283 & 22283 & 22283 & 1.0000 & 1.0000 & 1.0000 & 12.21 & 12.21 \\
062 & 1723314 & 1690862 & 1690862 & 1.0192 & 1.0000 & 1.0192 & 4.39 & 4.39 \\
063 & 612094 & 602954 & 602954 & 1.0152 & 1.0000 & 1.0152 & 5.80 & 5.80 \\
064 & 22504 & 22504 & 21944 & 1.0000 & 1.0255 & 1.0255 & 2.92 & 3.17 \\
065 & 11396 & 11396 & 11396 & 1.0000 & 1.0000 & 1.0000 & 6.64 & 6.64 \\
066 & 300629 & 300603 & 300603 & 1.0001 & 1.0000 & 1.0001 & 15.15 & 15.15 \\
067 & 15931687 & 15931687 & 15931687 & 1.0000 & 1.0000 & 1.0000 & 0.00 & 0.00 \\
068 & 362095 & 362095 & 362095 & 1.0000 & 1.0000 & 1.0000 & 0.00 & 0.00 \\
069 & 28700 & 28558 & 28540 & 1.0050 & 1.0006 & 1.0056 & 16.66 & 16.95 \\
070 & 26275 & 26216 & 26216 & 1.0023 & 1.0000 & 1.0023 & 28.70 & 28.70 \\
071 & 21551 & 21529 & 21529 & 1.0010 & 1.0000 & 1.0010 & 11.93 & 11.93 \\
072 & 4831281 & 4809827 & 4809827 & 1.0045 & 1.0000 & 1.0045 & 20.41 & 20.41 \\
073 & 3517515 & 3517515 & 3517515 & 1.0000 & 1.0000 & 1.0000 & 0.00 & 0.00 \\
074 & 337275 & 335626 & 335626 & 1.0049 & 1.0000 & 1.0049 & 31.17 & 31.17 \\
075 & 1430487 & 1430487 & 1430487 & 1.0000 & 1.0000 & 1.0000 & 0.03 & 0.03 \\
076 & 1240011 & 1232435 & 1232435 & 1.0061 & 1.0000 & 1.0061 & 44.61 & 44.61 \\
077 & 115617 & 115617 & 115617 & 1.0000 & 1.0000 & 1.0000 & 7.44 & 7.44 \\
078 & 1586167 & 1556908 & 1556908 & 1.0188 & 1.0000 & 1.0188 & 30.19 & 30.19 \\
079 & 6075976 & 6075472 & 6075472 & 1.0001 & 1.0000 & 1.0001 & 24.42 & 24.42 \\
080 & 113455 & 90408 & 90408 & 1.2549 & 1.0000 & 1.2549 & 34.41 & 34.41 \\
081 & 30943 & 30943 & 30943 & 1.0000 & 1.0000 & 1.0000 & 0.00 & 0.00 \\
082 & 721041 & 721041 & 721041 & 1.0000 & 1.0000 & 1.0000 & 0.00 & 0.00 \\
083 & 620731 & 620094 & 620094 & 1.0010 & 1.0000 & 1.0010 & 3.57 & 3.57 \\
084 & 503472 & 503472 & 503472 & 1.0000 & 1.0000 & 1.0000 & 0.00 & 0.00 \\
085 & 3282015 & 3281953 & 3281953 & 1.0000 & 1.0000 & 1.0000 & 11.65 & 11.65 \\
086 & 118403 & 118403 & 118403 & 1.0000 & 1.0000 & 1.0000 & 0.00 & 0.00 \\
087 & 833349 & 829669 & 829669 & 1.0044 & 1.0000 & 1.0044 & 30.22 & 30.22 \\
088 & 243623 & 243623 & 243623 & 1.0000 & 1.0000 & 1.0000 & 0.00 & 0.00 \\
089 & 147432 & 147432 & 147432 & 1.0000 & 1.0000 & 1.0000 & 0.00 & 0.00 \\
090 & 542638 & 509924 & 509924 & 1.0642 & 1.0000 & 1.0642 & 8.14 & 8.14 \\
091 & 1882479 & 1876653 & 1876653 & 1.0031 & 1.0000 & 1.0031 & 31.18 & 31.18 \\
092 & 1649430 & 1351088 & 1351088 & 1.2208 & 1.0000 & 1.2208 & 58.07 & 58.07 \\
093 & 15876 & 15876 & 15876 & 1.0000 & 1.0000 & 1.0000 & 0.00 & 0.00 \\
094 & 21573 & 20380 & 20380 & 1.0585 & 1.0000 & 1.0585 & 30.61 & 30.61 \\
095 & 420852 & 420852 & 420852 & 1.0000 & 1.0000 & 1.0000 & 0.00 & 0.00 \\
096 & 27053 & 27053 & 27053 & 1.0000 & 1.0000 & 1.0000 & 25.40 & 25.40 \\
097 & 2102348 & 2102250 & 2102250 & 1.0000 & 1.0000 & 1.0000 & 24.52 & 24.52 \\
098 & 90168 & 90168 & 90168 & 1.0000 & 1.0000 & 1.0000 & 0.00 & 0.00 \\
099 & 339475 & 329676 & 329676 & 1.0297 & 1.0000 & 1.0297 & 54.45 & 54.45 \\
100 & 58326 & 58085 & 58085 & 1.0041 & 1.0000 & 1.0041 & 33.19 & 33.19 \\
\end{longtable}

\endgroup

\section{Cache参数敏感性}\label{app:sensitivity}
六图各在四个非默认硬件配置下固定五核B最终计划，并尝试候选重优化。
下表的“失败候选”表示未通过依赖检查的记录；这些记录不赋予工期。
\begingroup
\small
\setlength{\tabcolsep}{4.5pt}
\renewcommand{\arraystretch}{0.94}
\setlength{\LTleft}{\fill}\setlength{\LTright}{\fill}
\begin{longtable}{crrrrr}
\caption{6张代表图的Cache参数敏感性}\label{tab:sensitivity-detail}\\
\toprule
图 & 参数 & 固定计划/周期 & 重优化/周期 & 固定命中/\% & 失败候选 \\
\midrule
\endfirsthead
\multicolumn{6}{c}{续表}\\
\toprule
图 & 参数 & 固定计划/周期 & 重优化/周期 & 固定命中/\% & 失败候选 \\
\midrule
\endhead
\midrule\multicolumn{6}{r}{续下页}\\
\endfoot
\bottomrule
\endlastfoot
012 & 带宽×2 & 13406 & 13406 & 7.68 & 0 \\
012 & 带宽/2 & 13406 & 13406 & 7.68 & 0 \\
012 & 容量×2 & 13406 & 13406 & 7.68 & 0 \\
012 & 容量/2 & 13406 & 13406 & 7.68 & 0 \\
016 & 带宽×2 & 7740714 & 7740714 & 33.33 & 1 \\
016 & 带宽/2 & 7740666 & 7740666 & 33.33 & 1 \\
016 & 容量×2 & 7740666 & 7740666 & 33.33 & 1 \\
016 & 容量/2 & 7740666 & 7740666 & 25.81 & 1 \\
051 & 带宽×2 & 627568 & 627568 & 42.39 & 1 \\
051 & 带宽/2 & 627371 & 627371 & 42.39 & 1 \\
051 & 容量×2 & 627502 & 627502 & 42.39 & 1 \\
051 & 容量/2 & 627502 & 627502 & 29.35 & 1 \\
062 & 带宽×2 & 1691572 & 1691572 & 4.39 & 1 \\
062 & 带宽/2 & 1691090 & 1691090 & 4.39 & 1 \\
062 & 容量×2 & 1685616 & 1685616 & 4.97 & 1 \\
062 & 容量/2 & 1718264 & 1718264 & 1.43 & 1 \\
067 & 带宽×2 & 15931687 & 15931687 & 0.00 & 0 \\
067 & 带宽/2 & 15931687 & 15931687 & 0.00 & 0 \\
067 & 容量×2 & 15931687 & 15931687 & 0.00 & 0 \\
067 & 容量/2 & 15931687 & 15931687 & 0.00 & 0 \\
081 & 带宽×2 & 30943 & 30943 & 0.00 & 0 \\
081 & 带宽/2 & 30943 & 30943 & 0.00 & 0 \\
081 & 容量×2 & 30943 & 30943 & 0.00 & 0 \\
081 & 容量/2 & 30943 & 30943 & 0.00 & 0 \\
\end{longtable}

\endgroup

\section{计算参数与配对字段}
正式评价使用题面默认的L1、UB、DDR和L2 Cache参数；敏感性实验仅改变L2容量或读带宽，其他计划和硬件条件保持不变。C初始命中率与C最终命中率分别保存，固定计划和重调度计划的字段由此分开。未通过依赖检查的候选不进入工期统计。
